% ============================================================================
% PAPER 9 — the philosophy-of-science extraction (FIRST-PERSON GENRE, by design)
% PROVENANCE: assembled by Claude Fable 5 (July 2026) at the author's commission
% from the author's own PoS content in the long manuscript (the criteria, the
% aphorisms, the genealogy, the journal entries) — author prose is the spine;
% external-case argumentation drafted; every thesis argued on historical cases
% BEFORE the author's own case per the landing blueprint agreed in-session.
% The dated audit travels in the companion session record.
% GENRE NOTE: unlike the other extractions this paper is deliberately first-
% person — the memoir-as-longitudinal-data method IS the contribution
% (lineage: Einstein's Autobiographical Notes, Poincare, Hadamard, Polanyi).
% R830 (Opus 5.5, Oct 21 2026): RESTRUCTURED after the external referee report (Direction A, auditable
% discovery records). See the [P9-R28 / R830] entry at the end of the round log. The pre-revision text is
% preserved verbatim, with a disposition note per section, in phil_sci_deferred.tex.
% R831 (Opus 5.5, Oct 21 2026): REVISED after the referee's second-round report (coding rules, the first
% transition's entailment table, atomic claims, frozen sources). See the [P9-R29 / R831] entry in the round log.
% ROUND LOG (P9-R3, this transcript; staged by the author's prep doc,
% drafting-executed; dispositions convention adopted per review): Item 1
% EXECUTED in the safe arbitration form — counts deferred to the provenance
% record, which arbitrates; the reconciled-specific alternative stands open
% as an author ruling. Item 2 EXECUTED (Spelke core-knowledge engagement
% installed at §8, adapted to the innateness frame per the prep doc's ADAPT
% license). Item 3 EXECUTED WITH A FLAGGED DEVIATION: the prep doc staged
% Science and Hypothesis (1905), but the invention essays named in the
% lineage live in Science and Method (1908), the entry already in the
% bibliography — cited accordingly, the 1905 entry not added; author
% override reserved. Item 4 EXECUTED (no candidate re-identification
% collectively ratified; follow-on sentence synced). Stack (a)–(g)
% EXECUTED, with (e) split per the prep doc's own condition: The Mind of
% the Ape observation cited to Premack & Premack 1983, the 1978 paper
% retained at the theory-of-mind construct in §8; Récoltes entry migrated
% to the Gallimard 2022 publication, body cite updated to the new key.
% CORRIGENDUM (drafting-side self-caught BEFORE this transcript's review;
% the original channel tag ``verification pass, this transcript,
% review-caught'' predated the review it credited --- relabeled in review
% per the corpus's false-provenance precedent): (i) the
% executed Item 4 carried an unstaged flourish beyond the agreed form ---
% ``the standing candidates ... are, notably, single-mind formulations
% themselves'' --- reverted: QBism (Fuchs--Caves--Schack) and GRW
% (Ghirardi--Rimini--Weber) are multi-author candidate re-identifications,
% so the flourish overclaims; the staged sentence (``no performed candidate
% ... collectively ratified'') is restored, author override reserved if a
% scoped version is wanted. (ii) This log's item-(e) line reads ``The Mind
% of the Ape''; the title is The Mind of an Ape, correct in the bibliography.
% [R-AUDIT, this transcript, review]: all P9-R3 items verified in place
% (installs, grammar, keys, residual hygiene all clean); corrigendum items
% (i) and (ii) CONFIRMED on the facts (QBism and GRW are indeed multi-author
% re-identifications; the 1983 title is The Mind of an Ape); the Poincare
% 1905-to-1908 deviation independently converged with the review's own
% catch; the corrigendum's channel label corrected as noted above.
% [P9-R4] (this transcript; author-commissioned, review-drafted --- a channel
% first for this paper, body prose drafted by the review side at the author's
% direct request, the author's approval of the wording the standing condition):
% the disclosure's syllogism paragraph installed --- the kinds-section
% completion (single-mind regularity + unratified move + the companions'
% claim) named and owned in the disclosure itself, its status priced exactly:
% a claim the companion papers must earn under independent refereeing, not a
% finding of this paper, whose theses stand or fall identically either way.
% [P9-R5] (this transcript; AUTHOR-ORIGINATED point --- the criteria-first,
% built-to-the-spec direction of dependence --- review-drafted wording with
% review-added consequence structure: the zero-evidence price stated and
% accepted, the priority claim graded per the paper's own dating standard
% [AUTHOR FILL: dated first statements of the criteria, if the record holds
% them, upgrade the grading clause to cited certificates], the two-stage
% exemption noted as undesignable; installed between the hazard disclosure
% and the syllogism paragraph; author approval of wording the standing
% condition per the P9-R4 precedent).
% [P9-R6] (this transcript; BRANCH RECONCILIATION + commissioned install,
% review-executed): the off-session branch's parallel execution of R3--R5
% arrived with variant wordings and is adopted as authoritative per the
% author's transmission (the Spelke rewrite --- object permanence, the
% metaphysical-reading sharpening; the section-9 tally rewritten around the
% branch-divergence fact, a sentence now SELF-INSTANTIATING, the two branches'
% divergent executions of these very rounds being an instance of the
% divergence it reports; the Grothendieck key retained at 1986 with the
% Gallimard note; the corroboration clause trimmed; the boundedness closure
% installed at the trig instantiation, its splice redundancy trimmed in
% adoption). REVIEW-CAUGHT in the parallel copy, NOT ADOPTED here: that
% copy's kinds-section follow-on reverted to ``unperformed by anyone''
% against its own Everett/Bohm clause --- this copy retains the resynced
% ``no performed candidate ... collectively ratified,'' and the parallel
% site should be repaired to match (second occurrence of this exact desync).
% DECLINED-WITH-RULING: the parallel copy silently discharged VERIFY-at-merge
% flags on ten external entries with no logged verification event --- flags
% retained here, re-dischargeable the moment the record logs who verified
% what; the Grothendieck note keeps its flag likewise. INSTALLED per the
% author's commission: the reflexivity-positioning footnote at the
% practices-both-halves sentence (Bloor's fourth tenet, Ashmore's performed
% thesis, the certificate-versus-regress delta), two bibliography entries
% added and flagged; review-drafted, author approval of wording the standing
% condition per precedent. TIMESTAMP RULING: the branch report's
% timestamp-ready verdict is accepted CONDITIONAL on this round --- freeze
% after P9-R6 lands, not before, a timestamp's office being to date a draft
% one can defend.
% [P9-R7] (this transcript; author-commissioned installs + a VERIFICATION
% EVENT, review-executed): (1) the 2023 print and its annotated edition made
% clickable at the framework's site URL; (2) the QBism on-ramp added to the
% second column --- the 2023 opening assumes QBism, the mature framework is
% non-doxastic, and the paper's own click-invitation made the omission a
% liability; (3) the correspondence archive installed as a third evidential
% leg, phrased precisely: the May 8, 2023 arXiv confirmation dates the
% readiness judgment, not the completion. VERIFICATION EVENT, logged: the
% section-7.1 claims about the 2023 print were checked this transcript
% against the deposited PDF fetched from the framework's site --- all three
% numbered postulates VERIFIED; the third's exists-forall two-generation
% alternation VERIFIED verbatim; the too-strong/too-weak elimination present
% ON THE PAGE; permanence flat, multiply; multi-observer consistency run
% through Postulate 3 explicitly; both superposition species extensionally
% present and insisted distinct (``to Aethic logic it is certainly
% distinct''); the second postulate's statistical-independence dress
% verbatim; Anti-Time present; the Rare-Earth application present WITH the
% Accordance Principle named in public print (a two-year-earlier public
% anchor); the 2023 modal lexical shadow also holds (zero counterfactual/
% Lewis/Stalnaker/perturbation). Remaining gate-side: the 2024 counts,
% the ``nowhere after'' halves, day-precision of May 20. Cross-reference:
% the annotated edition's body text was found abridged/smoothed relative to
% the deposited PDF and its verbatim restoration is performed in this same
% transcript (separate artifact; no Paper 9 claim depended on it).
% [P9-R8] (this transcript; author-commissioned): first-person source links
% installed so the record's artifacts are clickable at the claims that lean
% on them --- (1) the 2023 print and its annotated edition deep-linked
% directly (PDF + marked edition) in place of the hub URL; (2) the 2024
% version-1 claim now carries PhilArchive's public version trail, upgrading
% the successive-deposits claim to third-party-dated and clickable (the
% archive assigns the dates, not the author); (3) the correspondence
% archive's URL attached at its sentence; (4) the grade-school compilation
% linked at the memoir-grading clause WITH the grade preserved in the same
% breath (``memoir with a URL is still memoir'' --- the link must not
% inflate the evidence class); (5) the Aethus bib entry's note gains the
% version-trail URL. SEQUENCING NOTE for the timestamp: the annotated-
% edition link points at history/marked_original.html --- deploy the
% verbatim-restored edition (built this transcript) BEFORE freezing this
% file, so the linked artifact at freeze time is the restored one.
% [P9-R9] (this transcript; author-supplied provenance clarification,
% review-drafted at commission): the QBism second-column clause sharpened
% with direction-of-dependence. Author's clarification: Postulates 1--2
% developed in isolation in 2022; QBism and RQM read only in the months
% before the May 2023 write-up; extrusion-principle and soccer-field
% material already in hand August 2022. Grading applied strictly ---
% the half the record certifies (the dated Aug 22, 2022 entries predate
% the print by nine months, so the dress is provably younger than the
% content) is stated as record; the half only recollection supplies (the
% date of first interpretive-literature contact) is labeled ``the author
% reports.'' This closes the residual hostile reading the R7 clause left
% open (that QBism was the germ rather than the on-ramp). AUTHOR FILL
% (added to the standing list): any dated artifact of first QBism/RQM
% contact --- a journal mention, notes, a dated bookmark --- would upgrade
% the author's-report half to a certificate; supply if it exists.
% [P9-R10] (this transcript; author-disclosed artifact + VERIFICATION EVENT,
% review-executed at commission): the January 28, 2023 manuscript ``Centric
% Unfolding'' (hosted; fetched and read this transcript) installed as the
% record's fourth public print and the genealogy's missing intermediate.
% VERIFIED against the fetched PDF: the three rules of its Assembly section
% match VERBATIM, all three, the sentences the February 8 entry boxes as
% ``the old laws'' (journal-print interlock, now publicly checkable); the
% third rule is parent-directed (ancestral validity) --- the dated prior
% candidate of the elimination that yields the two-generation form May 8;
% the second rule runs through a reified Anti-Ethus complement DEFINED by
% statistical independence (since retired); soccer-field experiment in
% manuscript; a named compression postulate; permanence of the past Ethus;
% ZERO citations, ZERO references section, ZERO interpretive literature ---
% no QBism, no RQM; modal lexical shadow holds (zero counterfactual/Lewis/
% Stalnaker/possible-world). Consequences installed: public leg three->four
% prints (intro + checkability paragraph); January description added to the
% division-of-labor passage; QBism clause upgraded from author's-report to
% RECORD-BRACKETED (content dated on both sides of the dress's arrival) ---
% this DISCHARGES the R9 author-fill for a dated pre-May artifact, in the
% strong direction: not a dated first-contact, but a dated content-without-
% dress print, which is the better certificate; genealogy gains the Jan 28
% stop with the interlock stated; surprise-chain gains two links; May-vs-
% January disambiguation pass run. Date grading: month on the face, day and
% second (1/28/23, 9:51:29 AM) author-supplied file metadata --- same
% grading as May 20. Cross-reference: the aethic-code index card for this
% print was built in this same transcript. Gate-side: eyes-on confirmation
% that the file metadata reads as stated.
% [P9-R11] (this transcript; author-commissioned tone pass): the
% dare-the-world register removed while preserving the firewall's normative
% force. The ``merciless'' refrain softened uniformly at all four sites
% (abstract, §1 firewall statement, §6 division-of-labor note, closing
% prescription) to ``exacting'' --- uniform so the refrain remains a
% refrain; and ``cited, attacked, or used'' -> ``cited, challenged, or
% used.'' DELIBERATELY KEPT, judged metaphor rather than dare: the
% outlaw-and-sheriff closer; ``attack the alpha-space'' (idiomatic
% attack-a-problem usage, aimed at the terrain, not at readers); the
% heart-and-spine ``kills instantly'' anatomy. Rationale: an audit
% invitation should read as confidence in the record, not as a taunt to
% the auditor.
% [P9-R12] (this transcript; author-directed): (a) TITLE respecified under
% the author's separability requirement --- the philosophy must not read as
% inseparable from the framework; the paper is about the study of what
% makes inductive reasoning good, and the genealogy is its TEST CASE, not
% its subject. Main title retained (already fully general); subtitle now
% states the study and marks the dataset's role: ``On What Makes Inductive
% Reasoning Good --- Named Criteria, with a Dated, Auditable Genealogy as
% Test Case.'' ``Documented'' also retired for ``Dated, Auditable'' (the
% record's actual grade). (b) The substance/reasoning BOX restored to the
% long manuscript's concise two-sentence form VERBATIM (author-supplied),
% with the installed-to-rescue/entailed refinement demoted to a footnote
% via footnotemark/footnotetext (box-safe); the section's prose re-derives
% the direction-of-dependence formulation on the cases, so nothing is
% orphaned. Momentum ruling: boxes carry the quotable form; footnotes
% carry the precision.
% [R12 AMENDMENT, same transcript; author-rejected R12a]: the R12a title
% was a reviewer error twice over --- far too long for a title, and ``What
% Makes Inductive Reasoning Good'' both misstates the subject (the paper
% studies the structure of inference-driven discovery, not the goodness of
% inductive reasoning) and imports a colloquial register the face of a
% paper should not carry. Logged symmetrically. Installed instead: the
% author's original construction, trimmed --- main title unchanged;
% subtitle ``Criteria and a Dated Genealogy for Inference-Driven
% Discovery'' (five words shorter than the original subtitle, ``Dated''
% carrying the record claim, the subject named by the author's own
% original phrase).
% [R12 TITLE, FINAL --- author-supplied]: main title set to ``Inductive
% Reasoning as Art, Not Algorithm'' (the author's own phrase). Rationale
% on the record: it states the culminating art thesis in six words, names
% the exact contrast the Simon/BACON engagement argues, stays fully
% separable from the framework, and --- with the retained subtitle --- the
% full title enacts the split firewall: art for the making before the
% colon, criteria and a dated record for the checking after it. Read with
% §\ref{sec:art}'s machine-reframing exit in view: art-not-algorithm is a
% claim about the character of the reasoning, not about who or what can
% perform it.
% [P9-R13] (this transcript; author-originated, review-drafted at
% commission): the PROXIMITY-HARDNESS COROLLARY installed at the close of
% §shades --- short gaps are disproportionately hard to close because a
% dense neighborhood reads as coverage of the point; the firewall debate
% installed as the worked instance (standoff external history, closure
% this paper's; vacancy claim HEDGED ``so far as I have been able to
% find''); the evaluation discipline (vacancy-duration-under-density, not
% textual-delta size, measures difficulty) and the honesty discipline (the
% conditional runs one way: proximity-camouflage explains non-closure
% GIVEN correctness and never evidences correctness) both stated in the
% author's scoping (``a subset of applications''). Provenance of the
% exchange logged honestly: the reviewer's prior circularity objection
% targeted the unconditional form; the author's conditional refinement
% survives it and sharpens the thesis (shades now covers the last mile,
% not only the far distance); the one-way-conditional guardrail is
% retained in print at the reviewer's insistence, as the sentence that
% keeps the corollary out of the conventionalist stratagem's reach.
% [P9-R14] (this transcript; author-originated, review-drafted at
% commission): (a) the CO-SIGNATURE LOOP --- the dynamic form of the
% proximity corollary --- installed following it in §shades: fierce
% two-sided debate promotes the camps' shared-wrong alongside their
% shared-right (co-signature on the uncontested background); closure
% requires cutting possibility space exactly along the shared-right/
% shared-wrong boundary, excruciating from within and compounded by the
% loop; firewall instance deepened (the co-signed assumption was the
% coupling of evidential to normative status; the split position is the
% cut); the author's high-dimensional reading installed AS IMAGE per the
% paper's graph-conjecture precedent, closing on the observable (fierce-
% ness and vacancy co-occur). (b) the MASKING-DANCE FOOTNOTE installed in
% §9 on the review-loop clause: the R13 exchange narrated beat-by-beat as
% the corollary's first data point, citing the deposited session record
% for the verbatim; graded deliberately --- footnote placement per the
% author's own instinct (on the record, no main-body attention), the
% artificial-interlocutor caveat retained at the reviewer's insistence
% (extends scope, does not establish; exactly one instance). Rhetorical
% use disclaimed in the text itself.
% [R14 AMENDMENT, same transcript; author-commissioned]: the dance
% footnote expanded from narration to RECONSTRUCTION --- the misparse now
% articulated inline at the level of the words (the author's phrasing
% quoted verbatim, capitals preserved and marked; the surface chain the
% reviewer's extraction keyed on; the high-prior lexical slot it resolved
% to; the three sites where the real meaning was pressured: the
% antecedent's scope, the ``similar but different'' compression, the
% evaluative-vs-evidential register of ``count''); the bidirectional lock
% named per the Grothendieck exhibit; the deposited-record citation
% demoted from deferral to corroboration. Quote fidelity: fragments
% verbatim from the session exchange, ellipses none, capitals his.
% [P9-R15] (this transcript; author-approved, author-refined): the
% CLOSE-MEANS-LOW-SIGNIFICANCE fallacy warned against in the corollary's
% evaluation discipline, in the author's own coinage, WITH the author's
% anti-axiomatization point incorporated as the load-bearing reason the
% correction cannot be a sign-flip: a proximity-scoring rule would
% axiomatize a leg of inductive evaluation, violating the art thesis ---
% so the correction is an added COLUMN (delta statement + vacancy report:
% duration, density) with the unstated-vs-stated-and-rejected
% discriminator carrying the weight, and the §claim vacancy hedge made
% load-bearing (its defeat condition --- a produced prior statement ---
% announced in advance). Prescriptions gain the matching item, closing on
% ``the column is evidence, not algorithm.'' Reviewer note for the
% record: the discriminator supplies the corollary's cleanest falsifier
% and was reviewer-drafted; the refusal of the sign-flip and its
% art-thesis grounding are the author's.
% [P9-R16] (this transcript; author-originated articulation, review-
% drafted): WHY the sign-flip violates the art thesis, made exact by the
% author's monotonicity point --- a positive comparative principle is
% monotone, so affirming ``closer is harder'' asserts ``farther is
% easier'' in the same breath (a commitment nobody holds), while the
% negative form ``do not assume closer is easier'' blocks one fallacious
% inference and asserts no ordering; error-blocking negatives sit on the
% tips side of §art's heuristics-are-not-axioms line (the Vonnegut
% register --- passage grep-verified in-session before wiring to it).
% Consistency install alongside (reviewer-suggested, same license): the
% Prescriptions opener now claims the tips register explicitly, closing
% the latent ``no rules of discovery, yet a prescriptions section''
% gotcha before a referee opens it. Scope note: the asymmetry is stated
% for this instance, not as a taxonomy of all tips (Vonnegut's own rules
% include positives; the claim is that negatives of the error-blocking
% species are the genre at its safest, not that the genre is exclusively
% negative).
% [R16 AMENDMENT, same transcript; author-clarified]: the negative's
% content made exact and SYMMETRIC per the author --- the tip calls false
% BOTH universals (``closer is universally easier'' and ``farther is
% universally easier''), the pair the fallacy and its mirror respectively
% require, leaving the terrain order-free: structure removed, none
% installed. Reviewer-added and flagged as such: the recursion sentence
% (the two condemned universals co-sign the assumption that difficulty is
% monotone in proximity; the tip is the cut along that co-signature
% boundary --- the corollary's machinery run one level down on its own
% evaluation rule). Author's eyes requested on the recursion sentence at
% the reading pass: it is true and compact, but it is also the paper's
% third self-application in one section, and the author is the right
% judge of whether that density delights or fatigues.
% [P9-R17] (this transcript; author-originated exhibit, review-drafted):
% DELTA-EFFICACY installed in §criteria --- the Strong Implication
% delta's operational yield at first contact (the speedups-permitted
% inference convicting classical reductionism via the substitution test,
% during the introductory double-slit video, August 2022), tied to the
% vacancy column as its cash-out on this criterion specifically, with the
% §shades one-way conditional restated in place (yield dates effect;
% validity still rests on external cases). GRADING applied strictly:
% account attested in the 2024 deposit; TIMELINE INTERLOCK stated at
% two-hop precision after a reviewer near-miss caught pre-install (the
% draft's ``fortnight... exactly where within-the-week places them''
% conflated two chained week-steps; corrected to: video->eureka within
% the week per the deposit, eureka->entry within the week per the
% genealogy, twelve days total, hops closing cleanly); day-precision
% August 10 author-recollection, graded in text. Genealogy gains the
% first-contact stop (R17c) with the same grading. REVIEWER CAUTIONS,
% logged: ``thin-film'' not used as a Spekkens coinage (unverified as his
% term; the shape stated instead, spekkens2007evidence cited, VERIFY at
% merge); ``after thirty years'' softened to ``decades'' (figure not
% certifiable). MERGE RECOMMENDATION, not installed: the Sudoku analogy
% in §criteria carries \cite{spekkens2018quantum} in the long
% manuscript; migrate that citation and its bib entry at merge for
% attribution parity (not installed now to avoid a placeholder entry
% pre-freeze). AUTHOR FILL: any dated artifact of the Aug 10 viewing
% (watch history, a journal line) upgrades the day to certificate.
% [P9-R18] (this transcript; author-supplied): (a) the August 10 SOURCE
% disclosed --- read off the author's YouTube watch history in a lookup
% circa two years downstream --- upgrading the day's grade from bare
% recollection to platform-dated log with author-reported lookup; R17's
% AUTHOR FILL half-discharged (artifact identified), remainder narrowed:
% export or capture the history entry and hash it into the record.
% Installed at both sites (efficacy exhibit + genealogy stop, ``export
% pending''). (b) the NEXIC TRAINING conjecture installed as the
% exhibit's closing answer to how a first sitting arrives equipped:
% prior axiomatization practice training the meet-in-the-middle posture
% (working locally under an assumed global foundation; the author's
% phrasing ``gesturing at the outline of a foundation'' carried nearly
% verbatim); GRADED at the author's own hedge (``probably'') as
% introspective conjecture about one's own interior, Grothendieck-
% precedent invoked in the text itself; the certifiable half split out
% explicitly (Nexic practice priority, dated wing artifacts) from the
% conjectural half (the causal link). Content-before-name clause added,
% extending the record's signature from constructs to criteria.
% [R18 AMENDMENT, same transcript; author-directed]: (a) COURTESY reword
% --- ``the literature spent decades distilling'' -> ``took decades to
% distill,'' placing the duration on the field's timeline rather than on
% anyone's effort, per the author's explicit instruction not to read as
% a slight against Spekkens; (b) the ``thin film'' COINAGE restored as
% Spekkens' own, source author-supplied (Escaped Sapiens Podcast ep. 74,
% Nov 2024) and entered verbatim as spekkens2024escapedsapiens, cited
% alongside the toy-theory paper at the phrase; R17's reviewer caution
% on the coinage hereby DISCHARGED by author-supplied source --- the
% correct discharge direction: withheld pending verification, restored
% upon it. Entry carried without a VERIFY flag as author-supplied;
% spekkens2007evidence retains its VERIFY-at-merge flag as reviewer-
% added.
% [V1-FREEZE LEDGER] (this transcript): the author reports the timestamp
% OBTAINED on the R18-amendment state; that state is hereby the v1
% boundary. Gate-side at merge: record the frozen artifact's URL and
% hash on this line. All rounds from R19 forward are v2 content atop the
% frozen v1.
% [P9-R19] (v2; this transcript; author-commissioned literature loading
% + defense subsection): (a) STANFORD installed in §shades --- the
% unconceived-alternatives relation stated as division of labor with
% deltas at both ends (his record: THAT alternatives go unconceived and
% an anti-realist induction; ours: WHICH and WHY, plus an evaluation
% practice, no side taken on realism); bib stanford2006exceeding added,
% VERIFY at merge (reviewer-added). Includes the RETRODICTION CHECK ---
% unconceived alternatives should cluster where fierce debates crowded
% densest, testable on Stanford's own published cases --- v2's first
% dated out-of-sample commitment, on external material. (b) THE SHAPE OF
% THE DEFENSE subsection installed at §1's close: the calibrated prior
% conceded first; closure offered as the rebuttal with SIX interlocks
% listed, all verified on this paper's pages this transcript (the
% rooting-derives-implication toy derivation is NOT among them --- it
% lives only in the long manuscript; V2 CANDIDATE: migrate it into
% §criteria, which would add the strongest member to the list); the
% against-me consequence (correlation as shared risk, wounds travel,
% referee invited to test how far) stated before the for-me consequence
% (audit efficiency); one-way conditional closes.
% [P9-R20] (v2; author-originated thesis, review-drafted): the
% MEASUREMENT COROLLARY installed closing §art, framed as DERIVED from
% the art thesis (proactive reliable scoring of frontier capacity would
% span the making; scorer-that-reliable = generator-in-disguise = rules
% of discovery, denied), with the author's empirical anchor carried
% nearly verbatim (statistical evaluators are not emitting theories of
% everything -> metric fallibility entailed -> uncredentialed-outlier
% cases EXPECTED, credentials as retrospective proxy in proactive
% service, ``the metric caught in the act'' per the author's phrasing).
% Pre-empted objection built in: the verification asymmetry (P-vs-NP
% flavor) concerns DELIVERED results; the corollary targets UNDELIVERED
% ones (scoring requires modeling the prover). Person-level tie to the
% vacancy discount stated. CANDIDATE INSTANCE graded strictly: the 2023
% moderation decline (correspondence archive) named as the metric
% operating; false-negative status explicitly NOT claimed --- base rate
% only, adjudication assigned to the external record, one-way
% conditional in place. §shape gains the SEVENTH interlock (corollary
% derived-not-asserted). KIN logged as v2 citation candidates, not
% cited pending bib verification: Spence 1973 (signaling; proxies lossy
% by design), Collins 1979 (The Credential Society), Meehl 1954
% (clinical vs statistical prediction --- WITH the domain caveat that
% Meehl's superiority results live in criterion-ready tasks, not
% frontier discovery, so no conflict), Goodhart's law (adjacent, not
% identical: proxy corruption under targeting vs proxy insufficiency at
% the frontier).
% [P9-R21] (v2; author-commissioned queue discharge): (a) the ROOTING-
% DERIVES-IMPLICATION toy derivation MIGRATED from the long manuscript's
% critical-examination footnote into §criteria's close --- the author's
% own argument, adapted to this paper's register, with the toy-marking
% preserved and sharpened (informal quantifiers disclosed; the
% separability step's reliance on rooting's strongest reading disclosed);
% closes on ``one commitment viewed at two depths.'' R19's v2 candidate
% hereby DISCHARGED. (b) §shape gains its EIGHTH interlock, placed
% FIRST as the strongest relation on the list (derivation > policing).
% (c) the measurement corollary's KIN cited in-text with the Meehl
% domain caveat carried as the delta statement (criterion-ready tasks
% vs the unspanned frontier); four bib entries added, all VERIFY at
% merge, all reviewer-added --- the Goodhart source in its original 1975
% Reserve Bank of Australia form, flagged especially for verification
% since the law is usually cited from later reprints. QUEUE REMAINDER,
% off-paper by nature: actually RUNNING the Stanford retrodiction on his
% published cases --- an activity, not a text item; the paper's side
% (the advance statement) is complete.
% [P9-R22] (v2; author-commissioned: the queue's activity item executed):
% the Stanford retrodiction RUN and reported as a subsection closing
% section shades. VERIFICATION EVENT: Stanford's case series confirmed by
% web search this session against his own publications list --- Darwin
% pangenesis (BJPS 57), Galton stirp theory (Biology and Philosophy 21),
% Weismann germ-plasm (HPLS 27) --- and his proposed mechanism confirmed
% as COGNITIVE-GENERAL (failure to consider alternatives possibly a
% general property of human cognition), which is what makes the loop a
% structural amplifier rather than a rival and supplies the location-based
% discriminator the check needs. Corpus check also run: no companion
% paper engages Stanford, so R19 stands as the corpus's first engagement.
% RESULTS as installed: core series codes FIT (contested axis = soft vs
% hard heredity; co-signature = heredity/development as the parceling out
% of representative material; accepted alternative negates the
% co-signature); light, electrodynamics-of-moving-bodies and drift code
% FIT; heat codes PARTIAL and is reported as partial. UNBIDDEN FINDING:
% the same series sorts correctly on the substance/reasoning axis (failed
% theories add substance installed to rescue; the accepted alternative
% adds reasoning with its factors entailed) --- an external test on cases
% selected by someone else for another purpose. TWO CORRECTIONS AGAINST
% OURSELVES, installed rather than buried: (i) the R15 discriminator
% gains a third bin (stated-but-unengaged: Mendel 1866-1900, masking not
% rejection) and its second bin is narrowed (drift: correct rejection of
% the stated form, vacancy persisting for the modified version), with the
% R15 site amended to point here; (ii) the run is self-coded on partly
% familiar cases and is labeled a PILOT, not a result. PRE-REGISTERED
% PROTOCOL for an independent run installed in the text (fixed case list,
% no-stake coders, primary sources, stated miss condition), plus the
% requirement to engage the critical re-analysis of the same series.
% BIB: three Stanford case papers added (fields confirmed in-session,
% unflagged); novick2020presume added with AUTHORS flagged VERIFY --- the
% author list was not confirmed in-session and is deliberately omitted
% rather than guessed.
% [R22 AMENDMENT, same transcript; reviewer self-audit at author's
% instruction to verify the run lands]: four corrections to the run as
% first installed, three of them substantive. (1) FACTUAL: the claim that
% Stanford's book indexes the three additional cases was an overclaim ---
% the ether and caloric are among its background examples; light and
% continental drift are this paper's own additions and are now labeled
% as such. (2) CHRONOLOGICAL: the co-signature's defeat was compressed
% into one event; it arrived in two stages fifty years apart (segregation
% at the 1900 rediscovery, the retained-complement picture at
% mid-century), now stated, with the coding crediting the pair. (3)
% DENSITY: the soft/hard heredity axis was fiercely contested mainly in
% the 1880s-90s; at the Darwin end it sat nearer consensus, so the series
% supports the STATICS across its length but the DYNAMICS only in the
% later window --- installed as a qualification that cuts. (4)
% METHODOLOGICAL, the largest: the run tested the CONDITIONAL question
% (given a vacancy, was it in the co-signed complement?) and not the
% DISTRIBUTIONAL retrodiction actually stated in section shades
% (clustering), which needs a contrast class never assembled; and the run
% returned zero misses on cases selected for fame, so two untested
% candidate misses (phlogiston->oxygen, miasma->germ theory) are now
% named in the text for any independent run. Corrections count raised
% from two to four in the text. Reviewer note: all four were found by
% auditing the install rather than by external challenge, and the
% zero-miss return is the tell that prompted the audit.
% [P9-R23] (v2; author-supplied source + author-commissioned cross-
% referencing): the outlier-productivity paper (PhilPapers BENSAA-11,
% author-supplied code) read this session and cross-referenced at three
% sites. (a) REVIEWER ERROR CORRECTED, and it was a flattening of the
% author's position: R22's electrodynamics sentence read ``holding the
% transformations and much of the mathematics already,'' which silently
% adopted the Pais/Holton equal-footing framing --- ontology as the one
% brick still missing from an otherwise built wall --- the precise
% framing the author's Christmas-tree footnote rejects. Repaired: the
% standard historiography is now NAMED and CITED as the framing, the
% corpus's rejection installed with its image (the ontological substrate
% as the tree the mathematical pieces hang on; removal leaves broken
% glass, not a nearly finished theory) and its structural claim (the
% substrate IS the long sequential chain, not the short step left over),
% plus the author's ``in the air'' point --- an early landing manufactures
% the retrospective impression of near-completion --- and his closing
% formulation, that possessing every loose mathematical piece is at once
% why a problem looks solvable and why it stays unsolved. (b) The
% correction RESTRICTS this paper's own machinery and is installed as a
% restriction: the vacancy report must now declare the CURRENCY its
% density is counted in, since formal-apparatus proximity overstates
% closeness where the vacancy is ontological (R15 site amended). (c) The
% measurement corollary gains the interlock: capacity is measurable
% RETROSPECTIVELY (the study returns extreme values from achieved
% sequences) which is exactly why the PROACTIVE case is hopeless --- the
% retrospective score is computed from the output a proactive score would
% have to anticipate; no tension between the corollary and the study, and
% the pair is sharper than either alone. BIB: Outliers added in corpus
% self-citation format; pais1982subtle and holton1988thematic added with
% keys matched to the long manuscript's for merge parity, both flagged
% VERIFY at merge as reviewer-added.
% [P9-R24] (v2; external referee report, author-relayed; fixes executed
% where verifiable, flagged where they touch author-commissioned text).
% BRANCH-STATE NOTE first: this reviewer's model of the file was stale ---
% R21--R23 had landed from the drafting branch since R20, and the referee
% read the current state correctly. Verified before acting; no stale-base
% edits made. EXECUTED: (a) Q1, the abstract's ``three public prints'' ->
% ``four'' --- a REVIEWER ERROR of the R11 species, logged symmetrically:
% the R10 three->four sweep matched the intro's phrasing and missed the
% abstract's, a filtered-sweep miss now twice observed in this file; the
% abstract also given the body's grading language (third-party-dated
% where public, author-hosted metadata where self-deposited, memoir where
% recollection), closing the oversell the referee identified. (b) Q2, the
% historiographical contestability of the classic-case tellings conceded
% in \S1 with Newton's collectivity/gradualism named as the live rival
% reading, plus a footnote setting \S retrodiction's self-correcting
% standard as the bar the vignettes do not meet, and assigning the right
% inference (a rejected telling means the thesis owes a better case, not
% that it is refuted). (c) the referee's structural note: the pricing
% complex (corollary, co-signature loop, Stanford, retrodiction pilot)
% PROMOTED out of \S shades into its own section, so shades carries the
% masking thesis alone; six cross-references rerouted to sec:pricing,
% each checked individually (masking/lexicon/tree/orangutan refs left
% pointing at sec:shades). NOT EXECUTED, author decisions flagged in
% session: Q3's cut of the dance footnote (its expansion was author-
% commissioned precisely to avoid a deferral-promise; the referee's
% concern is length-as-signal, so an appendix relocation is the
% compromise offered) and the childhood compression.
% [R24 ADDENDUM, same transcript; referee's residual]: the closing
% provenance note's ``three-print continuity apparatus'' --- the THIRD
% sighting of the count-desync species this ledger has twice diagnosed,
% filed here per the referee's instruction. Disposition: DISAMBIGUATED,
% not updated --- the referent is the long manuscript's own apparatus,
% which genuinely tracked three prints at its deposit; the prose now
% names the count, the contrast with this record's four, and the reason
% (the apparatus predates the January print's surfacing), with the
% §record pointer attached. MERGE ITEM: update the long manuscript's
% continuity apparatus to four prints at merge. AUTHOR EYES-ON: confirm
% the three-print characterization of the long manuscript's apparatus
% is accurate as deposited.
% [P9-R25] (v2; author-corrected claim + author-supplied artifact, with
% a VERIFICATION EVENT this transcript): (a) ``years before either had a
% name'' -> ``months,'' per the author's correction that the two species
% were named by the turn of 2024 (May 20, 2023 -> ~7 months). The
% correction is not merely accepted, it is CORROBORATED at its midpoint:
% the October 19, 2023 video transcript, author-supplied and read this
% transcript, contains ZERO occurrences of ``agreeing'' or
% ``disagreeing'' while making the distinction extensionally --- so the
% unnamed interval is publicly attested at five months, with the naming
% date itself author-recollection and graded as such in the text.
% REVIEWER ERROR acknowledged: ``years'' was reviewer-drafted from the
% 2026 vantage without checking the naming date --- an over-claim in the
% author's favor, the species this ledger watches for hardest.
% (b) The October 19, 2023 PRESENTATIONS installed in \S record as an
% artifact of a DIFFERENT CLASS (explicitly not a fifth print; four
% prints stands), filling the public leg's largest gap (May 2023 ->
% Nov 2024). VERIFIED from the transcript: third postulate in finished
% two-generation form; second postulate still in statistical-
% independence dress; both species present, unnamed, distinction stated
% algebraically ((A and not-B) or (B and not-A) =/= (A and B)); speaker
% flags his own comprehension difficulty (masking's cost, self-reported
% in public, five months before naming); MODAL SHADOW EXTENDED TO
% SPEECH --- possible world / counterfactual / Lewis / Stalnaker /
% perturbation / many worlds / Copenhagen all zero. GRADING: video date
% platform-dated (YouTube), the strongest grade in the public leg after
% DKIM; the transcript is auto-generated captioning and visibly mangles
% the vocabulary (``ethis,'' ``ethys,'' ``arthus''), so NO verbatim
% quotation is taken from it and every claim above is stated in
% substance --- the caption artifact is why. AUTHOR EYES-ON: confirm
% October 19 is the premiere/publication date as displayed, and confirm
% the naming at the turn of 2024 (any dated artifact of first naming
% upgrades it to certificate). V2 CANDIDATES logged, not installed: the
% comprehension-difficulty line as a \S shades exhibit (self-reported
% masking cost, publicly dated); the Nexic video as the wings' earliest
% public co-presentation.
% [P9-R26] (v2; external referee round 2 --- Opus-family reviewer, noted:
% the review stack now spans two model families, a cross-model layer the
% provenance story gains for free). Five fixes + one residual, executed:
% (1) the stale ``below'' cured BY RELOCATION as prescribed --- the video
% paragraph moved to follow the 2024 lexical-shadow certificate, making
% the directional true without touching the sentence, restoring the
% four-print list's adjacency to ``the division of labor among the
% four,'' and pairing print-shadow with speech-shadow; (2) ``first
% public exposition'' -> ``first SPOKEN exposition'' (two public prints
% predate it; the contradiction was real); (3) the caption channel
% priced ON THE PAGE: counts declared as counts over auto-generated
% captions that garble coined vocabulary but not standard English ---
% the log's private grading promoted to text, closing the grade gap the
% referee named; (4) ``unscripted'' -> ``unscripted so far as the
% delivery indicates,'' and the date hedged ``as the platform displays
% it,'' matching the standing eyes-on item; (5) \S7.1 split into five
% paragraphs (four prescribed; the relocated video makes a natural
% fifth): prints / division of labor with certificates / the video
% artifact / second column / grading of the legs. Residual: the
% three-print disambiguation now dates the surfacing (``this record's
% 2026 assembly''). Referee's closing observation ledgered because it
% is the standard to hold: a new artifact is exactly where a corpus
% quietly upgrades its own evidence class --- R25--R26's gradings
% (different class, not a fifth print; channel priced; unscripted
% hedged) are this file's answer.
% [P9-R27] (v2; author-supplied passage with cleanup commission and an
% explicit layer-protection warning; layers inventoried BEFORE editing,
% per instruction, and recorded here). PROTECTED AND UNTOUCHED:
% ``unscrambling'' in the criterion (anticipates the schmid2020 citation
% so it lands as convergence); ``with no one counted'' (headcount pun);
% the exclusion obligation's absolutive payload landing last; the
% ratification move (openness defined sociologically --- no verdict
% observed --- so no interpretation is adjudicated); the Dirac-elegance
% concession; ``survives without fiat''; the scope-discipline double
% disclaimer; the footnote's register drop (rage disclosed, heat kept,
% ``babied their whole lives'' retained verbatim as voice). NEW MATERIAL
% INSTALLED (the paste's delta against the file): the irony sentence +
% the rage footnote, plus the paragraph split before ``Where I differ''
% present in the author's paste. CLEANUPS, each reversible: (E1) two
% grammatical ``as''es added to the classification sentence; (E2)
% ``speaks exclusively to'' -> ``speaks only to'' --- REVERSIBLE FLAG:
% the exclusion/exclusively echo may have been intentional; author to
% confirm or revert; (E3) the footnote's gerund chain (``the attempting
% to target... because of their reeling'') set into parseable clauses
% with content and heat preserved; main-text tail ``and then the same
% system doing so insists'' -> ``while the very system doing the
% muzzling insists'' (the contradiction now grips grammatically); (E4)
% the quantum sentence's NESTED dash pairs resolved by parenthesizing
% the Everett/Bohm aside --- content verbatim, bracket-matching
% restored. NOT changed though considered: ``borders on na\"ive''
% (author's cadence), ``when you also consider'' (the paper's direct
% address), the heart/spine anatomy (standing R11 disposition).
% [R27 AMENDMENT, same transcript; author-approved]: the footnote's
% closing clause scoped --- ``white men who were babied their whole
% lives'' -> ``the white men those same institutions babied'' --- per
% the calibration analysis: the universal was the one clause a bad-faith
% excerpt could lift, and the structural claim was always the intended
% one; the fix keeps ``babied'' and the heat, names the advantaged class
% within the institutions rather than as a blanket, and closes the
% double-bind into a single agent (the same institutions excluding and
% babying). ``Physical rage'' RETAINED by author decision after
% deliberation, with the epistemic-aptness-of-anger tradition (Lorde,
% Frye) noted in session as standing cover --- the register drop is
% disclosed interior at footnote grade, consistent with the corpus's
% grading discipline. Correction to the R27 log's framing, filed
% symmetrically: the reviewer described the paragraph split before
% ``Where I differ'' as new; the author clarifies the passage was
% already a paragraph he edited --- the split restores his structure
% rather than adds one.
% ============================================================================
% ============================================================================
% ============================================================================
% ============================================================================
% [R11 AMENDMENT, same transcript]: a FIFTH prose site found by the
% post-edit count check --- §9's ``verification was merciless'' --- which
% the initial sweep missed because its grep excluded comment lines and the
% sentence carries an inline attestation comment; softened to ``exacting''
% for refrain uniformity. Reviewer-side search-filter error, logged
% symmetrically per convention.
% ============================================================================
% ============================================================================
% [P9-R28 / R830] (Opus 5.5, Oct 21 2026; author-directed: "here's feedback from the ChatGPT referee about
% my phil_sci paper! Please push forward on revision that now"). RESTRUCTURED along Direction A of the report
% (auditable discovery records; the referee's preference; the ledger's earlier Paper-A blueprint). The reading
% was taken without a question to the author (resumed session) and stated to him; every cut block is preserved
% verbatim in phil_sci_deferred.tex, with the referee's notes per destination.
% TITLE: "Auditable Discovery Records: Telling Development from Reconstruction in a Dated First-Person Archive".
% ORIGINAL: "Inductive Reasoning as Art, Not Algorithm: Criteria and a Dated Genealogy for Inference-Driven
% Discovery". ABSTRACT rewritten last (<= 200 words, Perspectives on Science's limit).
% STRUCTURE: S1 one problem (development vs reconstruction), the contribution narrowed to the referee's
% combination, Dunbar / Polymath / Holmes / Medawar engaged, the disclosure compressed. S2 the qualified
% firewall only; the absolute version ("absent empirical consequences precommitted before the record was made,
% D does not license any change in P(H_T)") removed; predictivism (Maher; Hitchcock-Sober); the context
% distinction cited to Hoyningen-Huene, Schickore-Steinle, Aufrecht, and Norton, not claimed as a vacancy.
% S3 the G1-G4 grades REPLACED by an evidence-property table (after RFC 3161 and Haber-Stornetta) and an
% artifact register; "G1 and G2 establish priority against the world" WITHDRAWN; PhilArchive v1 an
% existence-by-date certificate (v1 14 Nov 2024 23:40:09 GMT; 519 versions to 19 Oct 2026, checked this round);
% four inferences tightened (absence of terms; no citations; the arXiv e-mail; later names). S4 a codebook fixed
% before coding. S5 three audited transitions with audit tables: the Jan->Feb 2023 rules (clarification), the
% Feb->May third postulate (substantive revision; the rejected one-generation candidate quoted), the two
% superpositions (renaming on the shared case, extension in the general definition); the second column extended.
% S6 six later accounts tested; one contradicted. S7 the two-stage model as a hypothesis with markers M1-M4 and
% alternatives; the record shows more than two stages; the immunity claim and the 2024->2026 out-of-sample claim
% WITHDRAWN. S8 the axis as a candidate coding distinction, Maxwell 1865 as the mixed case; masking as a coding
% hazard with the misunderstanding/disagreement test; the co-signature loop deferred with the referee's design
% requirements. S9 the art thesis bounded (four targets; the regress narrowed; Goedel machines; heuristics
% implementable; Norton). S10 limits, a methods statement with four statuses, supported practices only.
% Appendix A: coding sheet.
% SOURCES READ THIS ROUND: the January and May 2023 manuscripts as restored verbatim in the Marked Paper (the
% author's upload marked_original__2_.html, parsed: sheet-jan, sheet-may); the talk transcripts there (October
% 2023 talk at 39:30, 43:26, 46:42, 47:32; draft videos listed under June 12, 2024 at 15:34 and July 13, 2024 at
% 8:21); the Aug 22 2022, Oct 23 2022 and Feb 8 2023 entries as reproduced in the long manuscript.
% CORRECTIONS FOUND BY THE AUDIT: Jan 28 -> Feb 8 is ELEVEN days (the earlier draft and the Marked Paper's note
% say ten); the October 2023 talk is dated Oct 20 in the restoration's metadata and Oct 19 as displayed (time
% zone); the long manuscript's "the fundamental moment in which I first used the parent and child Aethic terms"
% (Feb 8) is contradicted by the January print's S4 (FLAG for the long manuscript; not edited this round).
% PROTECTED, KEPT: the first-person voice; the process/truth distinction; the built-to-spec disclosure and its
% zero-evidence price; the record's second column (extended); the exclusion paragraph and the RAGE FOOTNOTE
% (retained by author decision, R27), RELOCATED VERBATIM to S10.1, with a one-clause lead-in ("The notebook
% studies this paper sets itself beside are studies of celebrated individuals, and") and one closing sentence
% ("Whose records survive to be audited is itself a fact about access."); the outlaw/sheriff image (one line).
% DEFERRED (phil_sci_deferred.tex): the criteria (Strong Implication, Strong Rooting, the toy derivation); the
% axis's historical columns, graph conjecture and string theory; kinds of revolution (single-mind claim,
% inevitability, perceived advancedness); shades of similarity, branches, Grothendieck, apes, boundary
% prediction; proximity, vacancy, co-signature and the Stanford pilot; the art corollaries (alpha measure-zero,
% beauty, evaluator-generator) and the Hempel/Goodman paragraph; the memoir strata and the chalk-path
% through-line; entrenchment; the misparse footnote in full; the original prescriptions; the original abstract
% and S1.1. The R719 note (Oct 5 2026, "If Newton were here") had both candidate homes in deferred sections and
% stays pending with them.
% SECOND PASS (same day): an independent fresh-reader check verified every quotation against the restoration,
% the transcripts and the long manuscript, and found errors now fixed: Table 7 row 3 misdated ("blank" did not
% replace "statistically independent" by June 2024: the postulate is stated with "statistically independent" through
% July 2024 and the 2024 drafts DISTINGUISH the terms, "not statistically independent ... but it's still blank");
% the non-textbook reading of January's "statistically independent" now paired with the decomposition reading
% (January's components; the talk at 23:15) and row 5 recoded "untestable as stated"; the January print is the
% first MANUSCRIPT with three numbered rules (the October 2022 math-journal list is numbered too); the Feb 8 entry
% itself presents a rephrasing (the 2026 account inflates it); the detector predictions are not new in May
% (January derives both) and, by the talk's own "trial and error" (40:17), are accommodated --- stated in S5.2 as
% the one process fact bearing on truth; the two-stage markers scored on all five formulations (the agent drops
% out of the rule statements by October 2022; explicitness rises stepwise; the record fits ordinary learning
% with one larger step better than the hypothesis); the installed/constrained codes entered with competing
% readings (two of three contested); the earlier draft cited (PhilSciDraft) and its claims quoted exactly
% ("content before name"; the out-of-sample claim's conditional; the immunity claim's wording); Maxwell's
% "aethereal" per Phil. Trans.; "Nothing here bears on whether the framework is true" -> "does not assess".
% NOT DONE, FOR THE AUTHOR: independent human coding of the three transitions (protocol and sheet ready);
% conversion for Perspectives on Science (Word only; blind review, which a first-person record resists --- a
% point for the editor); the attestation items below.
% ============================================================================
% [P9-R29 / R831] (Opus 5.5, Oct 21 2026; author-directed: the referee's second-round report, "Here's the
% referee response, thanks!"). REVISED within Direction A along the report's seven points, its three priorities
% (coding definitions, the first transition's equivalence, frozen sources) first. The pre-R831 file is kept in the
% session backups; the deferred file is unchanged (it preserves the pre-R830 text).
% (1) CODES REPAIRED (S4, Table 3): component codes (preserved / renamed / added / dropped / indeterminate) and
% transition codes derived from them (renaming, clarification, extension, contraction, revision, mixed,
% abandonment). CLARIFICATION now requires TWO-WAY preservation (rule ii states the referee's A&B -> A case);
% MIXED (several determinate codes on different parts) is separated from INDETERMINATE (a component the texts do
% not settle); interpretive assumptions are entered with each code (rule iii) and competing readings are stated as
% the denial of a named assumption or a rival code for a named component (rule iv). Verdicts reduced to supported /
% contradicted / compatible, not discriminating / open, each with an evidence grade (independent / self-dated).
% TIMING: the rules were developed during this retrospective analysis and are frozen in this version, not
% preregistered; S1's "stated before they are applied" removed. S8.1: the Anti-Ethus code LEFT OPEN (not being
% anomaly-driven is not the positive condition; the case shows the distinction needs a rule for concept repairs);
% the compression postulate's August 21 2022 antecedent added ("I like the word `compress' better"); one code
% open, two contested.
% (2) THE FIRST TRANSITION (S5.1, Table 5): February components sought in January and January components sought
% in February. Each direction of law 2 is entered WITH ITS CONTRAPOSITIVE (one of the entry's two cases of
% invalidity), so logically equivalent rows can no longer be coded differently. Assumptions (FOUR): I1 (a contradiction in
% a parent's superposition makes the parent invalid; January's own "another way of putting it"), I2 (superposition =
% holding more than one value; "exists in all possible states at once", "which unlike the Anti-Ethus, is strictly
% forbidden"), I3, now a SCOPE assumption (January's rules state what holds of a realizable Ethus, so a breach makes
% an Ethus invalid rather than refuting the rule; January's treatment of rule 3 and footnote 2; no passage says it of
% rule 2, so I3 extends the pattern), I4 (footnote 3's check states a NECESSARY condition only; its "Therefore" follows
% from "an invalid Ethus ... may only beget invalid children Ethuses"; denied, the check is also sufficient, as the
% January detector-off derivation uses it, and the February laws, which state no sufficient condition for validity,
% drop it). New January rows: the S3 subdivision of partly correlated properties (carried by the term, rule i) and
% S4's "no contradictions are permitted" (preserved on I2; dropped if I2 is denied). Footnote 2 is
% mapped to law 2's converse through footnote 3 ("defeating its definition of independence"); the "never
% arbitrarily skew" clause is carried by the term (rule i; glossary); January's update mechanism ("reality will take
% the simplest possible path to avoid that, adding to the Ethus accordingly") is entered as OUTSIDE THE CONSTRUCT by
% a rule stated in S4 (an update statement is a different construct from a condition on states), with its track
% (Oct 2022 rule #5 -> January prose -> absent from the February laws -> May) and the note that counting it drops a
% component. Code: CLARIFICATION on I1-I4; CONTRACTION if I1, I3 or I4 is denied or the update rule is counted;
% REVISION if I2 is denied (the converse added, January's prohibition of contradiction dropped). The February entry's second reason ("doesn't
% work to hint at the clear correlation between the contradiction and the clash of Anti-Ethus superposition")
% quoted; the May "converse" passage entered as not settling I2 (it states the same direction twice).
% T2 EXTENSION, coded against the February laws, with TWO added components: the EA condition and validity by default
% ("If it also fails to prove your Aethus is invalid, then it must be valid"), relied on unstated in January on I4 and
% otherwise the sufficiency of footnote 3's check restored. The strongest competing reading denies S4's construct
% criterion: January's update rule arguably has the EA structure in dynamic form, restored in May; the T2 code
% stands (still added relative to February) but T1 becomes a contraction; the doomed-state criterion is new either
% way. May 8 is "an inference from the journal" (the May 7 entry already argues the detector case with the
% candidate); the January derivation of the detector results is from rules 2 and 3 with the uncertainty principle and
% footnote 3's check; the January rule is a derived consequence in May on I1.
% T3: MIXED if May left the "neither" case implicit (renaming on the shared case, extension in general), EXTENSION if
% the later texts add it; section retitled "a naming and a generalization"; the May
% detector-on passages read both ways are quoted; the Oct 2023 restatements in May's two-case wording (talk 48:26,
% summary) entered; the July 13 2024 drafts named (Solution to the Measurement Problem Part 1, 8:21; Part 2, 9:05
% detector on, 16:40 detector off). The second-column paragraph corrected: the Oct 2022 rule #5 left the numbered
% rules in January but its mechanism stayed in the prose and returns in May. HISTORICAL GLOSSARY (Table 4), with
% "Contradiction" in use from 2022 and January's own "basic laws" and "postulates".
% (3) LATER ACCOUNTS (S6, Table 8): six rows -> TWENTY-THREE ATOMIC CLAIMS. Eligibility stated (statements dating the
% audited constructs or describing at a date the artifacts in which they appear); the choice among eligible statements
% "was mine and followed no rule"; untested eligible statements named in the text and listed in the dossier (LM lines
% 10542, 10612, 10233, 10248, 10486, and recurrences of the May 8 dating). A WORDING column ("mine" / "drafted": the
% long manuscript's "For the Historian" (R145, Fable 5.1, Sept 2026), the annotation layer (Opus 4.6 draft, Opus 5
% cleaned, extended R159) and the earlier draft are drafted text adopted by the author; the genealogy list items and
% the author's notes are his) and evidence grades. Verdicts: 6 CONTRADICTED (three Feb 8 "firsts" in the author's
% wording; "ten days" (eleven); "as 'the old laws'" (the box is headed "Old Version"); the ED's May 20 "the same
% grading" as a date taken "by the file's metadata" (the PDF: 23 May 2023)), 7 SUPPORTED (two with notes), 4
% COMPATIBLE with named rivals, 6 OPEN (among them the first and third postulates "not changed since May 2023":
% unchanged in the Oct 2023 talk, restated by 2026, a transition not coded; and the May 8 dating, now compatible,
% not discriminating). Claim 1's evidence adds January's "basic laws" and "postulates"; claim 22's "concepts"
% corrected to "a 'concept'" (the August entry has the singular only).
% (4) TWO STAGES (S7): M4 is STABILITY (restatement as a principle in later versions; recurrence as a derived
% consequence not counted), not invariance under redescription (a different, unrun test), and on a track chosen to
% end at a surviving form it is fixed for the last formulation by the choice of track, so it carries no evidence
% there; M3 ANCHORED by definability (0 worked case only; 1 terms fixed by examples; 2 terms defined in the text or in
% an earlier text it cites). Table 9 rescored: Aug 2022 M2 absent, M3 2 (terms defined in the same paragraph); May 2023
% M3 1 ("contradict" fixed by example; January defines it but is not cited). The markers do not move together and
% explicitness does not rise steadily; the Oct 2022 row is now rule #1 with "that rule-breaking-portion ... must not
% exist" (rule #5's update mechanism is a different construct) and scores M4 YES (restated in January as "one final
% important rule of reality"); the stability clause is declared untestable on this track. The comparison with
% ordinary learning stays QUALITATIVE. The motive quoted from the contemporaneous Oct 23 2022 entry as well as the talk.
% Abstract (200 words): T1 "a clarification on four stated interpretive assumptions and a rule of scope, otherwise a
% contraction or a revision"; "the one track scored does not show the first". Verdict grades added to every compatible
% and open verdict; Table 8 group headings kept with their rows; version 1's content hedged "by my reading"; the LM
% caption states the v519 correspondence as the author's; the Einstein entry's month field fixed.
% (5) THE PAPER STANDS ON ITS OWN: references to the earlier draft removed except where it is tested evidence
% (S1 disclosure, S3 four inferences, S5.3, S7's surprise and out-of-sample paragraphs, S8.2, S9, S10.2); S9
% SHORTENED to two paragraphs, its box weakened to a description of the recorded practice (no claim that the art
% description beats procedural ones); the outlaw/sheriff line kept.
% (6) VERIFICATION: all quotations checked against the deposited 2023 PDFs, the journal export, the summary
% document, the restoration's transcripts, the long manuscript and the earlier draft (quotecheck_r831.py); the
% register gives the PDFs' metadata in UTC; the drafting ledger and provenance record linked (repository
% provenance/, commit 8e54a6e); fixed versions named (LM = PhilArchive v519, Oct 19 2026; AL = commit 8e54a6e; ED
% by SHA-256); the SUPPLEMENTARY SOURCE DOSSIER (phil_sci_dossier.md / .json / .pdf) freezes every quoted or tested
% passage with artifact, version, location and file hash. CORRECTIONS from the repository evidence: the
% Severed-Ethus is in the journal export (Jan 10 2023, 7:06 PM), not "known only from my later description";
% Anti-Time first appears in the Feb 19 2023 entry; the May 7 2023 entry states the one-generation condition as a
% "Postulate"; the May 9 2023 entry ties the arXiv account to "officially cracking the measurement problem"; the
% 2023 PDF URLs fixed (history/, not files/); the Sept 11 2026 hash manifest entered as identity, not dating.
% (7) SMALLER: S5.2 "the postulate accommodates the detector results ... no independent test", with Hitchcock and
% Sober on overfitting; S2 "confirms nothing" -> confirmatory force requires independent justification connecting
% the shape to a reliable process; S10.1 the archival-selection point now LEADS and CLOSES the exclusion paragraph;
% "ontology-capable individuals" GLOSSED in place from the author's own definition in the deferred kinds section
% (re-identifying what a theory's basic quantities are about) -- an insertion into the protected paragraph,
% flagged; the RAGE FOOTNOTE KEPT VERBATIM (author decision, R27); the referee's renewed objection (7d) is listed for
% the author below (the commented alternative is no longer in the source). The earlier framework dated 2018-2022.
% Length trimmed in S1, S4, S5.2, S7, S8 and S9 (still about 14,000 words with the tables the referee asked for).
% NOT DONE, FOR THE AUTHOR: the audio check of the coined names (July 13 2024, 8:21); version 1's term counts and
% names (PhilArchive's PDFs could not be fetched from the workspace; now stated "by my count"); confirmation that
% PhilArchive v519 is the Oct 19 deposit carrying the tested passages; deposit of the earlier draft and the session
% exports; testing the untested eligible statements (or confirming the stated selection); independent human
% coding; Word conversion and blinding for Perspectives on Science.
% ============================================================
% >>> EDIT (Fable 5, Aug 2026, author-directed): VERIFY-at-merge flags
% moved OUT of rendering bibliography note fields into this non-rendering
% ledger, so the PDF reference list carries no process tags. All flags
% remain in force for the merge pass.
% PURE FLAGS REMOVED (note field dropped), keys:
%   pais1982subtle --- VERIFY at merge; key matched to the long manuscript's for merge parity
%   holton1988thematic --- VERIFY at merge; key matched to the long manuscript's for merge parity
%   hempel1945studies --- VERIFY at merge
%   goodman1955fact --- VERIFY at merge
%   spence1973job --- VERIFY at merge
%   collins1979credential --- VERIFY at merge
%   meehl1954clinical --- VERIFY at merge
%   goodhart1975problems --- VERIFY at merge
%   novick2020presume --- AUTHORS VERIFY at merge --- title, journal, volume and number confirmed in-session; author list not confirmed and deliberately not guessed
%   stanford2006exceeding --- VERIFY at merge
%   spekkens2007evidence --- VERIFY at merge
%   vonnegut1999bagombo --- VERIFY at merge
%   lakatos1978methodology --- VERIFY at merge
%   hadamard1945invention --- VERIFY at merge
%   polanyi1958personal --- VERIFY at merge
%   kuhn1962structure --- VERIFY at merge
%   ashmore1989reflexive --- VERIFY at merge
%   schmid2020unscrambling --- VERIFY at merge; the corpus's recurring key family
%   reichenbach1938experience --- VERIFY at merge
%   popper1959logic --- VERIFY at merge
%   hanson1958patterns --- VERIFY at merge
%   gruber1974darwin --- VERIFY at merge
%   tweney1985faraday --- VERIFY at merge
%   holmes1991krebs --- VERIFY at merge
%   nersessian2008creating --- VERIFY at merge
%   nickles1980discovery --- VERIFY at merge
%   simon1973logic --- VERIFY at merge
%   langley1987scientific --- VERIFY at merge
%   chalmers2011verbal --- VERIFY at merge
%   deutsch2011beginning --- VERIFY at merge
%   duhem1954aim --- VERIFY at merge
%   hossenfelder2018lost --- VERIFY at merge
%   spelke2007core --- VERIFY at merge
%   premack1978does --- VERIFY at merge
%   premack1983mind --- VERIFY at merge
%   jordania2006who --- VERIFY at merge
%   jackson2004comme --- VERIFY at merge
%   poincare1908science --- VERIFY at merge
%   hacking1999social --- VERIFY at merge
%   feynman1960room --- VERIFY at merge
% COMPOSITE NOTES TRIMMED (substantive note text retained in PDF):
%   einstein1919times --- flag was: VERIFY at merge
%   bloor1976knowledge --- flag was: VERIFY at merge
%   grothendieck1986recoltes --- flag was: VERIFY at merge
%   jaynes1990probability --- flag was: VERIFY at merge
% R830 NEW KEYS (details checked this round where noted; the rest per standard references):
%   norton2003material; norton2021material (Calgary record checked); dunbar1995how (book title "The Nature of
%   Insight"; one secondary list prints "Mechanisms of insight"); dunbar1997how; gowers2009massively and
%   polymath2012new (checked); hitchcock2004prediction; hoyningen1987context; hoyningen2006context (Springer
%   page checked, pp. 119--131); schickore2006revisiting; aufrecht2017reichenbach (PhilArchive record checked);
%   holmes1987scientific (Isis 78, 220--235, doi 10.1086/354391); medawar1963fraud; haber1991timestamp;
%   rfc3161; rfc6376; schmidhuber2003godel (arXiv v5 checked); maxwell1865dynamical (Wikisource Part I quotes
%   checked); siegel1991innovation; MarkedPaper; Benander2023jan; Benander2023may; Benander2023talk.
%   maher1988prediction --- VERIFY pages (PhilPapers lists 272--285; 273--285 entered).
% R830 AUTHOR ATTESTATION: (1) the quotations of the 2023 manuscripts follow the 2026 verbatim restoration;
%   spot-check against the deposited PDFs; (2) version 1 carries the names agreeing/disagreeing superposition
%   and the two-generation third postulate (stated from the earlier draft; not re-read this round); (3) the
%   fifty-six / zero term counts in version 1 (carried from the earlier draft; recount); (4) the publicly
%   archived session report (completed) and the full record with hashes (planned), as S10.2 states them;
%   (5) the October 2023 talk's day (19 as displayed; 20 in the restoration's metadata); (6) S10.2 AUTHOR FILL:
%   the session report's URL; (7) S10.1 "ontology-capable individuals" lost its defining context with the kinds
%   section; kept verbatim as the author's protected paragraph --- keep, gloss, or reword; (8) the transcript
%   quotation of the coined names (13 July 2024, 8:21) to be checked against the audio; (9) PhilSciDraft --- the
%   September 2026 draft is a tested source (S3; Table 7, row 6); deposit it so the citation resolves publicly.
% R831 NEW KEYS: AethicCode (repository brunchwnewton/aethic-code, commit
%   8e54a6e978be3675d58acdf8485032d73818d297, Oct 20 2026; checked this round); Benander2023rules (the Google Doc
%   summary of Oct 20 2023, date per the site's listing; hosted export history/MajorRulesoftheAethus.html). The
%   Benander2023jan/may URLs moved to the commit-pinned repository paths (the /files/ paths do not exist in the
%   repository) with the PDFs' metadata dates; PhilSciDraft annotated with its SHA-256.
% R831 AUTHOR ATTESTATION: (1) the coined names at 8:21 of the July 13 2024 draft video, against the audio;
%   (2) version 1: the "possible world(s)" count (56), the zero counts, both superposition names, and the
%   two-generation third postulate; (3) PhilArchive v519 = the Oct 19 2026 long-manuscript source (SHA-256
%   21e24589357d7866...), or substitute the version that carries the passages tested in Table 8; (4) the gloss
%   inserted in S10.1 ("people able to re-identify what a theory's basic quantities are about"): keep, reword, or
%   remove; (5) the rage footnote, flagged again by the referee: keep (current) or delete the \footnote{...} in S10.1; (6) push the current long-manuscript source to the repository (the hosted copy predates Oct 19 and
%   lacks the note behind claims 21-22), and deposit the earlier draft and the session exports so the dossier's
%   hashes resolve publicly.
% ============================================================
\begin{filecontents}[overwrite]{phil_refs.bib}

@misc{Aethus, howpublished={Manuscript}, author={Benander, Ajax}, title={Aethic Reasoning: A Comprehensive Solution to the Quantum Measurement Problem}, year={2024}, month = {11}, url={https://philpapers.org/rec/BENARA-5}, eprint={https://philsci-archive.pitt.edu/id/eprint/25713}, eprinttype={Alternative URL}, note={Manuscript; the complete development of the framework, successive deposits November 2024--2026}}
@misc{AethusSemantics, author={Benander, Ajax}, title={Counterfactual Evaluation by Perturbation of Information States: Deleting the World-Similarity Ordering, Not Refining It}, year={2026}, howpublished={Preprint}, url={https://philpapers.org/rec/BENCEB-2}}
@misc{Outliers, author={Benander, Ajax}, title={Sequential Achievement and the Sharp Crossover in Outlier Productivity}, year={2026}, howpublished={Preprint}, url={https://philpapers.org/rec/BENSAA-11}}
@book{pais1982subtle, author = {Pais, Abraham}, title = {Subtle is the Lord: The Science and the Life of Albert Einstein}, publisher = {Oxford University Press}, address = {Oxford}, year = {1982}}
@book{holton1988thematic, author = {Holton, Gerald}, title = {Thematic Origins of Scientific Thought: Kepler to Einstein}, publisher = {Harvard University Press}, address = {Cambridge, MA}, edition = {Revised}, year = {1988}}
@misc{NexicFoundation, author={Benander, Ajax}, title={The Copernican Method Against the Copernican Principle: Anthropic Typicality over Partial Information States}, year={2026}, howpublished={Preprint}, url={https://philpapers.org/rec/BENTCM-7}}
@article{hempel1945studies, author = {Hempel, Carl G.}, title = {Studies in the Logic of Confirmation}, journal = {Mind}, volume = {54}, number = {213--214}, pages = {1--26, 97--121}, year = {1945}}
@book{goodman1955fact, author = {Goodman, Nelson}, title = {Fact, Fiction, and Forecast}, publisher = {Harvard University Press}, year = {1955}}
@article{spence1973job, author = {Spence, Michael}, title = {Job Market Signaling}, journal = {The Quarterly Journal of Economics}, volume = {87}, number = {3}, pages = {355--374}, year = {1973}}
@book{collins1979credential, author = {Collins, Randall}, title = {The Credential Society: An Historical Sociology of Education and Stratification}, publisher = {Academic Press}, address = {New York}, year = {1979}}
@book{meehl1954clinical, author = {Meehl, Paul E.}, title = {Clinical versus Statistical Prediction: A Theoretical Analysis and a Review of the Evidence}, publisher = {University of Minnesota Press}, address = {Minneapolis}, year = {1954}}
@incollection{goodhart1975problems, author = {Goodhart, Charles A. E.}, title = {Problems of Monetary Management: The {UK} Experience}, booktitle = {Papers in Monetary Economics}, volume = {I}, publisher = {Reserve Bank of Australia}, year = {1975}}
@article{stanford2006pangenesis, author = {Stanford, P. Kyle}, title = {Darwin's Pangenesis and the Problem of Unconceived Alternatives}, journal = {The British Journal for the Philosophy of Science}, volume = {57}, pages = {121--144}, year = {2006}}
@article{stanford2006galton, author = {Stanford, P. Kyle}, title = {Francis Galton's Stirp Theory of Inheritance and the Problem of Unconceived Alternatives}, journal = {Biology and Philosophy}, volume = {21}, pages = {523--536}, year = {2006}}
@article{stanford2005weismann, author = {Stanford, P. Kyle}, title = {August Weismann's Theory of the Germ-Plasm and the Problem of Unconceived Alternatives}, journal = {History and Philosophy of the Life Sciences}, volume = {27}, pages = {163--199}, year = {2005}}
@article{novick2020presume, title = {Presume It Not: True Causes in the Search for the Basis of Heredity}, journal = {The British Journal for the Philosophy of Science}, volume = {71}, number = {1}, pages = {59--86}, year = {2020}}
@book{stanford2006exceeding, author = {Stanford, P. Kyle}, title = {Exceeding Our Grasp: Science, History, and the Problem of Unconceived Alternatives}, publisher = {Oxford University Press}, address = {New York}, year = {2006}}
@misc{spekkens2024escapedsapiens, author = {Spekkens, Robert W.}, title = {What is Reality Made Of? {Unscrambling} the Quantum Omelet}, howpublished = {Interview by Shane Farnsworth, The Escaped Sapiens Podcast, Episode 74}, year = {2024}, month = nov, note = {Video available at \url{https://youtu.be/J2ZIRkfrFlI}}, url = {https://youtu.be/J2ZIRkfrFlI}}
@article{spekkens2007evidence, author = {Spekkens, Robert W.}, title = {Evidence for the epistemic view of quantum states: A toy theory}, journal = {Physical Review A}, volume = {75}, pages = {032110}, year = {2007}}
@book{vonnegut1999bagombo, author = {Vonnegut, Kurt}, title = {Bagombo Snuff Box}, publisher = {G.P. Putnam's Sons}, year = {1999}}
@article{einstein1919times, author = {Einstein, Albert}, title = {What is the Theory of Relativity?}, journal = {The Times (London)}, year = {1919}, month = {11}, note = {28 November 1919. The principle-theories versus constructive-theories distinction}}
@book{lakatos1978methodology, author = {Lakatos, Imre}, title = {The Methodology of Scientific Research Programmes}, publisher = {Cambridge University Press}, year = {1978}}
@book{hadamard1945invention, author = {Hadamard, Jacques}, title = {An Essay on the Psychology of Invention in the Mathematical Field}, publisher = {Princeton University Press}, year = {1945}}
@book{polanyi1958personal, author = {Polanyi, Michael}, title = {Personal Knowledge: Towards a Post-Critical Philosophy}, publisher = {University of Chicago Press}, year = {1958}}
@book{kuhn1962structure, author = {Kuhn, Thomas S.}, title = {The Structure of Scientific Revolutions}, publisher = {University of Chicago Press}, year = {1962}}
@book{bloor1976knowledge, author = {Bloor, David}, title = {Knowledge and Social Imagery}, publisher = {Routledge \& Kegan Paul}, year = {1976}, note = {The Strong Programme's reflexivity tenet}}
@book{ashmore1989reflexive, author = {Ashmore, Malcolm}, title = {The Reflexive Thesis: Wrighting Sociology of Scientific Knowledge}, publisher = {University of Chicago Press}, year = {1989}}
@article{schmid2020unscrambling, author = {Schmid, David and Selby, John H. and Spekkens, Robert W.}, title = {Unscrambling the omelette of causation and inference: The framework of causal-inferential theories}, journal = {arXiv:2009.03297}, year = {2020}}
@book{reichenbach1938experience, author = {Reichenbach, Hans}, title = {Experience and Prediction}, publisher = {University of Chicago Press}, year = {1938}}
@book{popper1959logic, author = {Popper, Karl R.}, title = {The Logic of Scientific Discovery}, publisher = {Hutchinson}, year = {1959}}
@book{hanson1958patterns, author = {Hanson, Norwood Russell}, title = {Patterns of Discovery}, publisher = {Cambridge University Press}, year = {1958}}
@book{gruber1974darwin, author = {Gruber, Howard E.}, title = {Darwin on Man: A Psychological Study of Scientific Creativity}, publisher = {E.P. Dutton}, year = {1974}}
@incollection{tweney1985faraday, author = {Tweney, Ryan D.}, title = {Faraday's Discovery of Induction: A Cognitive Approach}, booktitle = {Faraday Rediscovered}, publisher = {Macmillan}, year = {1985}}
@book{holmes1991krebs, author = {Holmes, Frederic Lawrence}, title = {Hans Krebs: The Formation of a Scientific Life}, publisher = {Oxford University Press}, year = {1991}}
@book{nersessian2008creating, author = {Nersessian, Nancy J.}, title = {Creating Scientific Concepts}, publisher = {MIT Press}, year = {2008}}
@book{nickles1980discovery, author = {Nickles, Thomas (ed.)}, title = {Scientific Discovery, Logic, and Rationality}, publisher = {Reidel}, year = {1980}}
@article{simon1973logic, author = {Simon, Herbert A.}, title = {Does Scientific Discovery Have a Logic?}, journal = {Philosophy of Science}, volume = {40}, number = {4}, pages = {471--480}, year = {1973}}
@book{langley1987scientific, author = {Langley, Pat and Simon, Herbert A. and Bradshaw, Gary L. and Zytkow, Jan M.}, title = {Scientific Discovery: Computational Explorations of the Creative Processes}, publisher = {MIT Press}, year = {1987}}
@article{chalmers2011verbal, author = {Chalmers, David J.}, title = {Verbal Disputes}, journal = {Philosophical Review}, volume = {120}, number = {4}, pages = {515--566}, year = {2011}}
@book{deutsch2011beginning, author = {Deutsch, David}, title = {The Beginning of Infinity}, publisher = {Allen Lane}, year = {2011}}
@book{duhem1954aim, author = {Duhem, Pierre}, title = {The Aim and Structure of Physical Theory}, publisher = {Princeton University Press}, year = {1954}}
@book{hossenfelder2018lost, author = {Hossenfelder, Sabine}, title = {Lost in Math: How Beauty Leads Physics Astray}, publisher = {Basic Books}, year = {2018}}
@article{spelke2007core, author = {Spelke, Elizabeth S. and Kinzler, Katherine D.}, title = {Core Knowledge}, journal = {Developmental Science}, volume = {10}, number = {1}, pages = {89--96}, year = {2007}}
@article{premack1978does, author = {Premack, David and Woodruff, Guy}, title = {Does the Chimpanzee Have a Theory of Mind?}, journal = {Behavioral and Brain Sciences}, volume = {1}, number = {4}, pages = {515--526}, year = {1978}}
@book{premack1983mind, author = {Premack, David and Premack, Ann James}, title = {The Mind of an Ape}, publisher = {W. W. Norton}, year = {1983}}
@book{jordania2006who, author = {Jordania, Joseph}, title = {Who Asked the First Question?}, publisher = {Logos}, year = {2006}}
@article{jackson2004comme, author = {Jackson, Allyn}, title = {Comme Appel\'e du N\'eant --- As If Summoned from the Void: The Life of Alexandre Grothendieck}, journal = {Notices of the AMS}, volume = {51}, number = {9--10}, year = {2004}}
@book{grothendieck1986recoltes, author = {Grothendieck, Alexander}, title = {R\'ecoltes et Semailles}, year = {1986}, note = {Circulated 1986; published Gallimard, 2022}}
@book{poincare1908science, author = {Poincar\'e, Henri}, title = {Science and Method}, publisher = {Thomas Nelson}, year = {1908}}
@article{jaynes1990probability, author = {Jaynes, E. T.}, title = {Probability in Quantum Theory}, year = {1990}, note = {In Complexity, Entropy, and the Physics of Information (W. H. Zurek, ed.), Addison-Wesley; the omelette passage}}
@book{hacking1999social, author = {Hacking, Ian}, title = {The Social Construction of What?}, publisher = {Harvard University Press}, year = {1999}}
@article{feynman1960room, author = {Feynman, Richard P.}, title = {There's Plenty of Room at the Bottom}, journal = {Engineering and Science}, volume = {23}, number = {5}, pages = {22--36}, year = {1960}}

@misc{AethusQuant,
  author       = {Benander, Ajax},
  title        = {Quantum Mechanics from Records: Three Consistency Postulates, One Composition Primitive, and One Kernel Recover the Measurement Architecture of the Dirac--von Neumann Formalism, with the Projection Postulate as a Theorem},
  year         = {2026},
  howpublished = {Preprint},
  url          = {https://philpapers.org/rec/BENTQB},
  note         = {The physics-face axiomatization; first posted July 17, 2026.}
}


@article{norton2003material, author = {Norton, John D.}, title = {A Material Theory of Induction}, journal = {Philosophy of Science}, volume = {70}, number = {4}, pages = {647--670}, year = {2003}}
@book{norton2021material, author = {Norton, John D.}, title = {The Material Theory of Induction}, publisher = {University of Calgary Press}, address = {Calgary}, series = {BSPS Open}, year = {2021}}
@incollection{dunbar1995how, author = {Dunbar, Kevin}, title = {How Scientists Really Reason: Scientific Reasoning in Real-World Laboratories}, booktitle = {The Nature of Insight}, editor = {Sternberg, Robert J. and Davidson, Janet E.}, publisher = {MIT Press}, address = {Cambridge, MA}, pages = {365--395}, year = {1995}}
@incollection{dunbar1997how, author = {Dunbar, Kevin}, title = {How Scientists Think: On-Line Creativity and Conceptual Change in Science}, booktitle = {Creative Thought: An Investigation of Conceptual Structures and Processes}, editor = {Ward, Thomas B. and Smith, Steven M. and Vaid, Jyotsna}, publisher = {American Psychological Association}, address = {Washington, DC}, pages = {461--493}, year = {1997}}
@article{gowers2009massively, author = {Gowers, Timothy and Nielsen, Michael}, title = {Massively Collaborative Mathematics}, journal = {Nature}, volume = {461}, number = {7266}, pages = {879--881}, year = {2009}}
@article{polymath2012new, author = {Polymath, D. H. J.}, title = {A New Proof of the Density {Hales--Jewett} Theorem}, journal = {Annals of Mathematics}, volume = {175}, number = {3}, pages = {1283--1327}, year = {2012}}
@inproceedings{maher1988prediction, author = {Maher, Patrick}, title = {Prediction, Accommodation, and the Logic of Discovery}, booktitle = {PSA 1988: Proceedings of the Biennial Meeting of the Philosophy of Science Association}, volume = {1}, pages = {273--285}, year = {1988}}
@article{hitchcock2004prediction, author = {Hitchcock, Christopher and Sober, Elliott}, title = {Prediction versus Accommodation and the Risk of Overfitting}, journal = {The British Journal for the Philosophy of Science}, volume = {55}, number = {1}, pages = {1--34}, year = {2004}}
@article{hoyningen1987context, author = {Hoyningen-Huene, Paul}, title = {Context of Discovery and Context of Justification}, journal = {Studies in History and Philosophy of Science}, volume = {18}, number = {4}, pages = {501--515}, year = {1987}}
@incollection{hoyningen2006context, author = {Hoyningen-Huene, Paul}, title = {Context of Discovery versus Context of Justification and {Thomas Kuhn}}, booktitle = {Revisiting Discovery and Justification: Historical and Philosophical Perspectives on the Context Distinction}, editor = {Schickore, Jutta and Steinle, Friedrich}, series = {Archimedes}, number = {14}, publisher = {Springer}, address = {Dordrecht}, pages = {119--131}, year = {2006}}
@book{schickore2006revisiting, editor = {Schickore, Jutta and Steinle, Friedrich}, title = {Revisiting Discovery and Justification: Historical and Philosophical Perspectives on the Context Distinction}, series = {Archimedes}, number = {14}, publisher = {Springer}, address = {Dordrecht}, year = {2006}}
@article{aufrecht2017reichenbach, author = {Aufrecht, Monica}, title = {Reichenbach Falls---and Rises? {Reconstructing} the Discovery\slash Justification Distinction}, journal = {International Studies in the Philosophy of Science}, volume = {31}, number = {2}, pages = {151--176}, year = {2017}}
@article{holmes1987scientific, author = {Holmes, Frederic L.}, title = {Scientific Writing and Scientific Discovery}, journal = {Isis}, volume = {78}, number = {2}, pages = {220--235}, year = {1987}}
@article{medawar1963fraud, author = {Medawar, Peter B.}, title = {Is the Scientific Paper a Fraud?}, journal = {The Listener}, volume = {70}, pages = {377--378}, year = {1963}}
@article{haber1991timestamp, author = {Haber, Stuart and Stornetta, W. Scott}, title = {How to Time-Stamp a Digital Document}, journal = {Journal of Cryptology}, volume = {3}, number = {2}, pages = {99--111}, year = {1991}}
@misc{rfc3161, author = {Adams, Carlisle and Cain, Pat and Pinkas, Denis and Zuccherato, Robert}, title = {Internet {X.509} Public Key Infrastructure Time-Stamp Protocol ({TSP})}, howpublished = {RFC 3161, Internet Engineering Task Force}, year = {2001}, url = {https://www.rfc-editor.org/rfc/rfc3161}}
@misc{rfc6376, author = {Crocker, Dave and Hansen, Tony and Kucherawy, Murray}, title = {{DomainKeys Identified Mail (DKIM)} Signatures}, howpublished = {RFC 6376, Internet Engineering Task Force}, year = {2011}, url = {https://www.rfc-editor.org/rfc/rfc6376}}
@misc{schmidhuber2003godel, author = {Schmidhuber, J{\"u}rgen}, title = {G{\"o}del Machines: Self-Referential Universal Problem Solvers Making Provably Optimal Self-Improvements}, howpublished = {arXiv:cs/0309048}, year = {2003}, note = {Version 5, December 2006}}
@article{maxwell1865dynamical, author = {Maxwell, James Clerk}, title = {A Dynamical Theory of the Electromagnetic Field}, journal = {Philosophical Transactions of the Royal Society of London}, volume = {155}, pages = {459--512}, year = {1865}}
@book{siegel1991innovation, author = {Siegel, Daniel M.}, title = {Innovation in {Maxwell's} Electromagnetic Theory: Molecular Vortices, Displacement Current, and Light}, publisher = {Cambridge University Press}, address = {Cambridge}, year = {1991}}
@misc{MarkedPaper, author = {Benander, Ajax}, title = {The Marked Paper: Linguistic Masking in {Ajax's} 2023 Manuscript}, howpublished = {Web edition: the January and May 2023 manuscripts restored verbatim, with a 2026 annotation layer and transcripts of talks}, year = {2026}, url = {https://aethus.org/aethic-code/history/marked_original.html}}
@misc{Benander2023jan, author = {Benander, Ajax}, title = {Centric Unfolding}, howpublished = {Manuscript}, year = {2023}, month = {1}, note = {PDF metadata 28 January 2023, 14:51:29 UTC; repository \texttt{history/Jan2023Aethic.pdf}}, url = {https://github.com/brunchwnewton/aethic-code/blob/8e54a6e978be3675d58acdf8485032d73818d297/history/Jan2023Aethic.pdf}}
@misc{Benander2023may, author = {Benander, Ajax}, title = {Solving the Measurement Problem: A Comprehensive Model Based on Three Postulates}, howpublished = {Manuscript}, year = {2023}, month = {5}, note = {PDF metadata 23 May 2023, 17:10:04 UTC; repository \texttt{history/Original\_Aethic.pdf}}, url = {https://github.com/brunchwnewton/aethic-code/blob/8e54a6e978be3675d58acdf8485032d73818d297/history/Original_Aethic.pdf}}
@misc{PhilSciDraft, author = {Benander, Ajax}, title = {Inductive Reasoning as Art, Not Algorithm: Criteria and a Dated Genealogy for Inference-Driven Discovery}, howpublished = {Unpublished draft}, year = {2026}, month = {9}, note = {Source file SHA-256 e9559f4a798a3f37\ldots; see the supplementary source dossier}}
@misc{Benander2023talk, author = {Benander, Ajax}, title = {Did {I} Just Solve the Measurement Problem?}, howpublished = {Video lecture, YouTube}, year = {2023}, month = {10}, url = {https://www.youtube.com/watch?v=_XDr8mNfGRk}}
@misc{AethicCode, author = {Benander, Ajax}, title = {Aethic Reasoning: Code and Demos}, howpublished = {GitHub repository, commit 8e54a6e978be3675d58acdf8485032d73818d297 (20 October 2026)}, year = {2026}, url = {https://github.com/brunchwnewton/aethic-code}}
@misc{Benander2023rules, author = {Benander, Ajax}, title = {Major Rules of the Aethus}, howpublished = {Google document, dated 20 October 2023 by the author's site; HTML export in the aethic-code repository, \texttt{history/MajorRulesoftheAethus.html}}, year = {2023}, month = {10}}
\end{filecontents}

\documentclass[11pt]{article}
\usepackage[margin=1.05in]{geometry}
\usepackage{amsmath, amssymb, amsthm}
\usepackage{booktabs, tabularx, array}
\usepackage{xltabular}
\setlength{\LTcapwidth}{\textwidth}
\usepackage{enumitem}
\usepackage[hidelinks]{hyperref}
\hypersetup{hypertexnames=false} % R831: unique anchors (longtables spanning pages duplicated table.N)
\usepackage{xurl}
\usepackage{microtype}
\IfFileExists{lmodern.sty}{\usepackage{lmodern}}{\microtypesetup{expansion=false}}
\usepackage{tcolorbox}
\usepackage[style=numeric, sorting=none]{biblatex}
\addbibresource{phil_refs.bib} % always loaded: PoS-specific keys live only here; sibling keys duplicated for standalone compiles carry MIGRATE notes for the merge

\newtcolorbox{criterion}[1]{colback=white, colframe=black, title={#1}, fonttitle=\bfseries}
\setlength{\emergencystretch}{2em}
\appto{\bibsetup}{\raggedright}


\title{Auditable Discovery Records:\\ \large Telling Development from Reconstruction in a Dated First-Person Archive}
\author{Ajax Benander\thanks{Email: \href{mailto:akb9408@rit.edu}{akb9408@rit.edu}, Website: \url{https://aethus.org}. This paper draws its data from the dated record of one framework's construction, much of it reproduced in the long companion manuscript \cite{Aethus}. It is written in the first person because its data are first-person (\S\ref{sec:claim}). Every passage it quotes or tests is frozen, with its location and the SHA-256 hash of its source file, in a supplementary source dossier.}}
\date{First drafted July 2026\\ {\small This version: October 2026}}

\begin{document}
\maketitle

\begin{abstract}
\noindent Records kept while a theory is being built are evidence about how it was built, and can test its builder's later account. This paper develops a protocol for such tests and applies it to my record of one framework's construction, from timestamped journal entries of 2022 to a preprint independently dated November 2024. The protocol separates what each kind of artifact can establish, codes changes between formulations component by component, and tests later accounts claim by claim. Three transitions are audited: a redesign of the framework's rules, a clarification on four stated interpretive assumptions and a rule of scope, otherwise a contraction or a revision; a change to its third postulate that adds a condition and rejects a too-strong candidate on the page; and a distinction drawn without a name in 2023, renamed and generalized in 2024. Of twenty-three claims selected from my later accounts, the record contradicts six, including three misdated firsts in my own wording. A two-stage account of breakthroughs and a distinction between installed and independently constrained structure are restated as a hypothesis and a coding aid, and the one track scored does not show the first. The paper does not assess whether the framework is true.
\end{abstract}

% ============================================================================
\section{The problem, and a disclosure}\label{sec:claim}

The discovery side of science is not unstudied. Its notebooks have been read closely and well: Gruber on Darwin \cite{gruber1974darwin}, Tweney on Faraday \cite{tweney1985faraday}, Holmes on Krebs \cite{holmes1991krebs}, Nersessian on Maxwell \cite{nersessian2008creating}, with Hanson and the discovery-logic debates collected by Nickles behind them \cite{hanson1958patterns, nickles1980discovery}. Reasoning has also been recorded as it happened. Dunbar spent a year in the weekly meetings of molecular-biology and immunology laboratories, studying how researchers used analogies and handled unexpected results \cite{dunbar1995how, dunbar1997how}. And mathematics now has public records of discovery in progress, written by the participants: the Polymath projects worked in timestamped blog comments, dead ends included, so that, as Gowers and Nielsen put it, ``for the first time one can see on full display a complete account of how a serious mathematical result was discovered'' \cite{gowers2009massively, polymath2012new}.

All of these records meet the same question when a later account is set beside them. Published papers present finished structures, and the complaint that they misrepresent the work behind them is old \cite{medawar1963fraud}; Holmes, reading Lavoisier's successive drafts, found new ideas arising in the act of writing \cite{holmes1987scientific}. Memoirs and retrospectives do the same thing at longer range. When a later formulation is set beside an earlier one, three things can be going on: the later text renames or clarifies what the earlier one already contained; the content changed; or the later account reads into the earlier text what was not there. This paper asks how a dated record can tell these apart, and answers with a protocol applied to one record.

The record is mine. It documents the construction, between 2022 and 2026, of a framework in the foundations of quantum theory \cite{Aethus, AethusQuant}: journal entries timestamped as they were written, two manuscripts of 2023, recorded talks, and a preprint deposited in November 2024 whose version history is kept by the archive rather than by me. A first-person record adds something third-party reconstruction lacks, the builder's own reading of his notes, and that reading can be tested against the notes themselves. It also carries an obvious hazard: the person auditing the record is the person whose story it tells. The design answers the hazard in three ways. Claims about the record are routed to artifacts a reader can inspect, and the passages used are frozen in a supplementary dossier; each code is entered with the interpretive assumptions it needs, so that a reader who disagrees can locate the disagreement; and the audit is run against my own later accounts, several of which it finds contradicted.

The contribution is accordingly narrow. Longitudinal first-person records are not new, nor are graded provenance and the worry about retrospective reconstruction. What this paper offers is their combination: a long first-person record, a separation of the properties each artifact can certify, explicit interpretive rules, developed during this analysis and frozen in this version, and tests of later accounts against earlier artifacts. Against the notebook studies, the record adds its author's own reading of his notes, open to test against them, and, for part of its span, dating by a third party. Against Dunbar's laboratory data, it adds a span of years and the private stages that come before any group meeting. Against Polymath, it adds an individual's pre-public stages, at the price that its private leg is dated by its author rather than by a platform. It has one subject, self-selected, and its codes are so far mine.

Now the disclosure, which the design makes mandatory rather than polite. I hold methodological views that classify my own framework favorably, and the framework was built to satisfy two criteria of adequacy that I held before it existed. This paper neither states nor uses those criteria, which I develop separately. A construction engineered to pass a test contributes nothing to the test's validity, so the framework's scores on those criteria are no evidence for them, and nothing here depends on them. This paper's conclusions concern process, what the record shows about how the framework was built, and a reader who rejects the framework should be able to assess them on the record alone. Whether the framework is true is a question for its own papers, under independent refereeing.

One declaration of genre. The paper is written in the first person because its data are first-person. The method has a small lineage --- Einstein's autobiographical notes, Poincar\'e's and Hadamard's essays on invention, Polanyi's personal knowledge \cite{poincare1908science, hadamard1945invention, polanyi1958personal} --- whose common understanding is that some evidence about invention exists only in the first person.\footnote{Reflexivity is not new to the study of science. The Strong Programme made it a tenet \cite{bloor1976knowledge}, and Ashmore performed it on the page \cite{ashmore1989reflexive}, where textual self-reference spiraled toward regress. The self-application here is of a different kind: it ends in comparisons between dated artifacts rather than in further text about itself.} Sections~\ref{sec:evidence}--\ref{sec:protocol} say what a discovery record can be evidence for, describe this one, and state the protocol; Section~\ref{sec:transitions} audits three transitions and Section~\ref{sec:reconstruction} tests later accounts; Sections~\ref{sec:stages}--\ref{sec:art} restate two of my earlier theses as a hypothesis and a coding aid and discuss judgment at points of revision; Section~\ref{sec:limits} gives the record's limits, a statement on AI assistance, and the practices the analysis supports.

% ============================================================================
\section{What a discovery record can be evidence for}\label{sec:evidence}

Let $D$ be discovery-record evidence, $H_P$ a hypothesis about the process by which a theory arose, and $H_T$ the hypothesis that the theory is true. $D$ bears directly on $H_P$: on what was written and when, on whether idea $A$ preceded idea $B$, on whether a later account matches the contemporaneous notes. Its bearing on $H_T$ is indirect. A record can show that a theory was fitted before a held-out dataset was seen, that an apparent derivation was reverse-engineered from its answer, that independent routes converged on one result, or that reporting was selective. Each is a fact about the process that rationally moves confidence in $H_T$, because each bears on the reliability of the method that produced the theory; the debate over prediction and accommodation turns on exactly such facts \cite{maher1988prediction, hitchcock2004prediction}.

Literal Bayesian screening, $P(H_T \mid D, E, M) = P(H_T \mid E, M)$, holds when $H_T$ is independent of $D$ given $(E, M)$, with $E$ the object-level evidence and $M$ independently validated information about the method's reliability. Where $D$ informs $M$, the record bears on truth through that channel. The shape of the story is another matter. A discovery narrative with a favored shape --- an epiphany, a surprise, an early insight later vindicated --- does not confirm the theory by having that shape. Whatever confirmatory force the shape has requires independent justification connecting it to a reliable process, and where that justification exists, the weight rests on the process the shape indicates rather than on the shape. Every question this paper puts to its record is a question about $H_P$, which is why the disclosure of Section~\ref{sec:claim} leaves the paper's conclusions standing, although one process fact the audit turns up bears on truth through the channel just described (Section~\ref{sec:t2}).

This position on the old separation of the context of discovery from the context of justification \cite{reichenbach1938experience, popper1959logic} drops the separation for content, since records of discovery are evidence about process, and keeps it for normativity, since a record's shape is no warrant by itself. The position is not new. Hoyningen-Huene, cataloguing the distinction's versions, located its core usefulness in separating descriptive from normative perspectives on science \cite{hoyningen1987context, hoyningen2006context}; the essays collected by Schickore and Steinle revisit it from historical and philosophical sides \cite{schickore2006revisiting}; and Aufrecht has argued that Reichenbach's own version does not map onto Hoyningen-Huene's, so that the distinction has no single clear meaning \cite{aufrecht2017reichenbach}. The nearest relative of the view that rejecting universal rules of induction need not mean abandoning rational assessment is Norton's material theory, on which there are no universal inductive schemas and the warrant of an induction lies in facts of its domain \cite{norton2003material, norton2021material}. Section~\ref{sec:art} says what the record adds to that view, and what it does not.

% ============================================================================
\section{The record, and what each artifact establishes}\label{sec:record}

The record has a private leg and a public leg. The private leg is a run of journal entries, several timestamped to the minute as they were written; the entries used here are in a public export of the journal in my repository \cite{AethicCode}, and most are also reproduced in the long companion manuscript \cite{Aethus}. The public leg begins with two manuscripts of 2023 \cite{Benander2023jan, Benander2023may}, hosted with a verbatim restoration and a separately marked annotation layer of 2026 \cite{MarkedPaper}. It continues with a summary document and recorded talks of October 2023 and after \cite{Benander2023rules, Benander2023talk}, and with a preprint deposited at PhilArchive on November 14, 2024, whose version trail records every later deposit, 519 of them by October 19, 2026 (\url{https://philarchive.org/versions/BENARA-5}). The record also reaches back into an earlier framework, built between 2018 and 2022, and, at memoir grade, into childhood; neither is audited here.

The artifacts differ in what they can establish, and the differences are easy to blur. Integrity (these are the bytes that were deposited), existence by a date, publication, the order of events, the presence or absence of words, and an author's later account are separate properties, and no single artifact certifies all of them. Timestamping protocols draw the line exactly: identifying data by a hash is one thing, and establishing that the data existed before a given time is another, which requires a trusted third party or an independently dated anchor \cite{haber1991timestamp, rfc3161}. Table~\ref{tab:evidence} separates the properties, with the record's instances, and Table~\ref{tab:register} lists the artifacts the audit uses.

\begin{table}[htbp]
\centering\small
\caption{Kinds of evidence in a discovery record, what each can support, and the record's instances.}\label{tab:evidence}
\begin{tabularx}{\textwidth}{@{}>{\raggedright\arraybackslash}p{0.17\textwidth}>{\raggedright\arraybackslash}X>{\raggedright\arraybackslash}X>{\raggedright\arraybackslash}p{0.23\textwidth}@{}}
\toprule
Evidence & What it can support & What it does not establish by itself & Instance in this record\\
\midrule
Independently dated deposit & That the deposited content existed by that date & First conception, originality, correctness & Preprint, version 1, PhilArchive, 14 Nov.\ 2024\\
Platform-dated upload & That the recording existed by the upload date & When its content was first conceived; a transcript's fidelity & Talks posted October 2023\\
Cryptographic hash & Identity and integrity of particular bytes & When the bytes were created, or when what they record happened & Hashes in the source dossier; a manifest hashing the math-journal master file, 11 Sept.\ 2026\\
Hash anchored in an independently dated record & That matching bytes existed by the anchor's date & That the record is complete & None for the private leg\\
Signed message from a third party (DKIM) & Origin domain and integrity; a date asserted by the sender's systems \cite{rfc6376} & What the event meant to its recipient & arXiv account confirmation, 8 May 2023\\
Local metadata or self-hosted dating & The chronology the author recorded & An independently certified chronology & Journal timestamps; PDF metadata of the 2023 manuscripts\\
Presence or absence of terms & That the terms occur, or do not, in the inspected text & Exposure to, or isolation from, the associated ideas & Term counts in version 1\\
Retrospective recollection & The author's later account of an episode & A contemporaneous record of the episode & The long manuscript's genealogy; the 2026 annotations; parts of this paper\\
\bottomrule
\end{tabularx}
\end{table}

\begin{table}[htbp]
\centering\small
\caption{The artifacts the audit uses. ``Self'' marks dating under the author's control.}\label{tab:register}
\begin{tabularx}{\textwidth}{@{}>{\raggedright\arraybackslash}p{0.29\textwidth}>{\raggedright\arraybackslash}p{0.2\textwidth}>{\raggedright\arraybackslash}p{0.13\textwidth}>{\raggedright\arraybackslash}X@{}}
\toprule
Artifact & Date & Dating basis & Used for\\
\midrule
Journal entry, ``Explanation For Double Slit Experiment with Centric Unfolding'' & 21 Aug.\ 2022, 10:27 AM & Self (journal) & ``Compress'' preferred to ``collapse''\\
Journal entry, ``Double Slit Experiment With Human Subjects'' & 22 Aug.\ 2022, 1:23 and 1:42 PM & Self (journal) & The unknown set dropped; the Anti-Ethus defined; the question of future states\\
Journal entry, ``Rules for Superposition'' & 23 Oct.\ 2022, 10:37 AM & Self (journal) & Rules \#1--\#5; the worry about communicating with the future\\
Journal entry, ``Severed-Ethus'' & 10 Jan.\ 2023, 7:06 PM & Self (journal) & A construct not carried forward\\
Manuscript, \emph{Centric Unfolding} & ``January 2023'' on its face; 28 Jan., 14:51:29 UTC by PDF metadata & Self (hosted) & Three rules; child and parent defined; the paradox of footnote~6\\
Journal entry, ``Rephrasing the Three Laws of the Ethus'' & 8 Feb.\ 2023, 4:42 PM & Self (journal) & Old and new laws, with reasons\\
Journal entry, ``Revisiting the Measurement Problem'' & 7 May 2023, 10:54 PM & Self (journal) & A one-generation candidate stated as a ``Postulate''\\
Account-confirmation message, arXiv & 8 May 2023 & Third party (DKIM) & An account event\\
Journal entry, ``Plan for arXiv Publish'' & 9 May 2023, 6:24 PM & Self (journal) & The account event tied to a judgment that the work was done\\
Manuscript, \emph{Solving the Measurement Problem} & ``May 2023'' on its face; 23 May, 17:10:04 UTC by PDF metadata & Self (hosted) & Three postulates; a rejected candidate; two superpositions\\
Summary document of the rules (Google Docs export) & 20 Oct.\ 2023 by the site's listing & Self (hosted export) & The superpositions unnamed; the ``neither'' case kept\\
Recorded talk, ``Did I Just Solve the Measurement Problem?'' & 19--20 Oct.\ 2023 (the displayed day varies with time zone) & Platform & The postulates in their May wording; the superpositions unnamed\\
Draft video series & May--Aug.\ 2024 & Self (listed on the author's site; platform record not linked) & ``Blank''; the names of the two superpositions\\
Preprint, version 1 & 14 Nov.\ 2024, 23:40:09 GMT & Archive & Existence by date; term counts\\
Long manuscript, version 519; annotation layer; earlier draft of this paper & 2026 & Archive; self; self & The later accounts tested in Section~\ref{sec:reconstruction}\\
\bottomrule
\end{tabularx}
\end{table}

The strongest independent anchor is the preprint's version trail. It is an existence-by-date certificate: it shows that the content of version~1 existed by November 14, 2024. Any stronger claim, that the content was mine first or first anywhere, needs separate argument, and this paper does not make one. Before the preprint, the independently dated artifacts are few: the talks posted in October and November 2023, and the signed messages of a correspondence archive, which date events such as the opening of an account rather than the content of the work. The journal entries and the two 2023 manuscripts are dated by their author, and their order is checkable for internal consistency and no more. Some of that checking can be done: the February 2023 entry quotes the January manuscript's three rules verbatim as its ``Old Version'' (Section~\ref{sec:t1}), which is what the January manuscript's priority predicts, and the entry of May 7, 2023 states as a ``Postulate'' the one-generation condition that the May manuscript proposes and rejects (Section~\ref{sec:t2}). A hash fixes bytes, not the dates of what the bytes record: the September 2026 manifest identifies the math-journal master file from the date its hash was recorded, and the entries inside the file remain dated by their author.

Four inferences from such artifacts are easily overdrawn. First, across version~1 of the preprint the phrase ``possible world(s)'' occurs, by my count, fifty-six times, while ``counterfactual,'' ``Lewis,'' ``Stalnaker,'' and ``perturbation'' do not occur (the dossier does not hold version~1, and lists these counts as unverified). That establishes what the text contains. It is compatible both with isolation from the modal literature and with exposure that went uncited, so it cannot establish isolation by itself. Second, the January 2023 manuscript has no citations; that establishes what it cites, not what its author had read. Third, a DKIM-signed message from arXiv dates the opening of an account on May 8, 2023. A journal entry dated the next day ties the event to a judgment that the work was done: ``yesterday, right after, (at least I hope), officially cracking the measurement problem, I made an arXiv account.'' The entry is self-dated, and the message dates only the event. Fourth, a later name for an earlier construction does not establish that the earlier and later contents are the same. That has to be shown, which Section~\ref{sec:t3} attempts for one pair of names.

Two derived artifacts need grading of their own. The talks are quoted from transcripts made by speech recognition and repaired by a logged rule list restricted to the framework's coined terms. The transcripts are a reading aid, and a quotation that matters should be checked against the audio at the timestamp given; below, only phrases in standard English are quoted from them, no coined term is counted from them, and the one exception, the first recorded use of two coined names, is flagged where it occurs. The 2023 manuscripts are quoted from the deposited PDFs, against which every quotation here has been checked; the verbatim restoration reproduces them, slips included, and keeps the 2026 annotations visibly apart from the text. The passages this paper quotes or tests are frozen in the supplementary dossier, each with its artifact, version, location, and the SHA-256 hash of its source file, so that a checker need not rely on any later edition.

% ============================================================================
\section{An audit protocol}\label{sec:protocol}

A transition is an ordered pair of formulations of one construct, each tied to an artifact in Table~\ref{tab:register}. Constructs are identified by function, not by name: the third rule of January 2023 and the third postulate of May 2023 are formulations of one construct because each is the framework's condition on which states can be realized. By the same criterion, a statement of how a state is updated belongs to a different construct from a statement of which states are real, in superposition, or valid. Each formulation is divided into components, claims that can be located in its text, and the coding proceeds component by component. A transition's code is derived from its components' codes, and a later account is tested claim by claim. Table~\ref{tab:codes} defines the codes and verdicts.

\begin{table}[htbp]
\centering\small
\caption{Codes for components and transitions, and verdicts for later accounts.}\label{tab:codes}
\begin{tabularx}{\textwidth}{@{}>{\raggedright\arraybackslash}p{0.2\textwidth}>{\raggedright\arraybackslash}X@{}}
\toprule
\multicolumn{2}{@{}l}{\emph{Component codes, for each component of either formulation}}\\
\midrule
Preserved & Stated or entailed in the other formulation, under the interpretation entered with the code\\
Renamed & Preserved under a new term whose extension, read by use (rule~i), is unchanged\\
Added & A component of the later formulation, neither stated nor entailed in the earlier\\
Dropped & A component of the earlier formulation, neither stated nor entailed in the later\\
Indeterminate & The texts do not settle which of the codes above applies; the rival codes are named\\
\midrule
\multicolumn{2}{@{}l}{\emph{Transition codes, derived from the component codes}}\\
\midrule
Renaming & Every component preserved, some under new terms, nothing else changed\\
Clarification & Every component of each formulation preserved in the other; wording, explicitness, or organization changed\\
Extension & Some component added, none dropped\\
Contraction & Some component dropped, none added\\
Revision & Components both added and dropped\\
Mixed & Parts of the transition, such as different cases, receive different codes; each part is reported with its own\\
Abandonment & A construct dropped without successor; where its work is absorbed, if anywhere, is recorded\\
\midrule
\multicolumn{2}{@{}l}{\emph{Verdicts, for each atomic claim of a later account}}\\
\midrule
Supported & An earlier artifact states or entails the claim\\
Contradicted & An earlier artifact states or entails its negation\\
Compatible, not discriminating & The earlier artifacts fit the claim and a named rival equally\\
Open & The record does not bear on the claim, or bears on it only through an unsettled code\\
\bottomrule
\multicolumn{2}{@{}p{\textwidth}@{}}{\footnotesize A code that holds only on an interpretive assumption is conditional on it. Each verdict carries the grade of its best evidence: independent, when a third party dates it, or self-dated.}\\
\end{tabularx}
\end{table}

Five rules govern the coding.
\begin{enumerate}[label=(\roman*), leftmargin=*, itemsep=2pt, topsep=4pt]
\item \emph{Sense by use.} A term's content is fixed by how the artifact uses it, not by its textbook sense and not by the sense a later text gives it. Table~\ref{tab:glossary} records the senses used here, and Section~\ref{sec:masking} explains why this rule carries weight in this record.
\item \emph{Two-way preservation.} A transition is a clarification only if every component of each formulation is preserved in the other. One-way entailment is not enough: a later formulation entailed by the earlier may have dropped some of its content, as $A$ is entailed by $A \land B$.
\item \emph{Stated interpretations.} Each component code is entered with the interpretive assumptions it needs. A code that holds only on an assumption is conditional on it, and the transition is reported with the code it receives if the assumption is denied.
\item \emph{Paired readings.} Each code is entered with the strongest competing reading, stated as the denial of a named assumption or as a rival code for a named component, so that a coder who disagrees can say where. A transition whose parts receive different determinate codes is mixed; a component whose code the texts do not settle is indeterminate. The first is a finding about the transition, the second about the evidence.
\item \emph{Outcome-blind coding, and direction of test.} A transition is coded from its two artifacts before consulting whether the construct survived; survival is recorded separately. A later account is split into atomic claims, each asserting one thing, and each claim is tested only against artifacts dated before the account.
\end{enumerate}

The codes and rules were developed during this retrospective analysis, by the person whose record it is, and are frozen in this version. They were not preregistered; their place before the transitions they are applied to is an order of exposition. The codes in Sections~\ref{sec:transitions} and~\ref{sec:reconstruction} cannot meet the first clause of rule~(v), since I chose the transitions knowing their outcomes, and they are offered as a first coding for an independent coder to test. A coder with no stake in the framework, working from the frozen passages with Tables~\ref{tab:codes} and~\ref{tab:glossary} and the coding sheet of Appendix~\ref{app:coding}, can recode the same transitions blind to my codes. Such a recoding tests the reliability of the codes, not the hypotheses of Sections~\ref{sec:stages} and~\ref{sec:hazards}, which need records other than mine.

% ============================================================================
\section{Three transitions}\label{sec:transitions}

The three transitions were chosen, with their outcomes known, to exercise the codes: a redesign whose code turns on stated interpretive assumptions, an extension with a rejected candidate on the page, and a naming whose code turns on one component the texts leave undetermined. All three concern the framework's postulates. The first two fall between January and May 2023, so their order rests on the private leg; the first independent anchor for any of their content is the October 2023 talk. Table~\ref{tab:glossary} glosses the framework's terms as the audited versions use them, so that the early texts are not read in the mature framework's senses.

\begin{table}[htbp]
\centering\small
\caption{The framework's terms as used in the versions audited.}\label{tab:glossary}
\begin{tabularx}{\textwidth}{@{}>{\raggedright\arraybackslash}p{0.16\textwidth}>{\raggedright\arraybackslash}p{0.15\textwidth}>{\raggedright\arraybackslash}X@{}}
\toprule
Term & In use & Sense in use, with the passage that fixes it\\
\midrule
Ethus & Aug.\ 2022 -- Feb.\ 2023 & A state of knowledge: ``Let this set of knowledge be called an Ethus'' (Jan.\ 2023); ``the abstract set of knowns something may hold on the world'' (journal, 10 Jan.\ 2023)\\
Aethus & May 2023 on & The same construct renamed: ``Let this reality-perceiving entity be referred to as an `Aethus' '' (May 2023)\\
Anti-Ethus & Aug.\ 2022 -- Jan.\ 2023 & ``the set of all possible things which are Rebetary, (Relatively Arbitrary / statistically independent), to the Ethus'' (Aug.\ 2022); it ``contains all the properties in reality which would have zero net effect on the Ethus if they were altered'' (Jan.\ 2023). Retired from the laws in Feb.\ 2023\\
Statistically independent & Aug.\ 2022 -- July 2024 & Glossed as ``Relatively Arbitrary'' (Aug.\ 2022); a property only partly correlated with the Ethus ``may simply be subdivided into an Ethus component and an Anti-Ethus component'' (Jan.\ 2023). Whether the term keeps its textbook sense is open (Section~\ref{sec:reconstruction})\\
Superposition & 2022 on & Holding every admissible value at once, skewed only as the state of knowledge directs: an Anti-Ethus variable ``exists in all possible states at once,'' and the superposition of the Anti-Ethus ``will never arbitrarily skew one way or another without the direct order of the Ethus to do so'' (Jan.\ 2023); an independent attribute ``will occupy all possible states at once that it may occupy without influencing the Aethus'' (May 2023)\\
Contradiction & 2022 on & Used from 2022 (``No contradictions allowed!'', Oct.\ 2022); defined in January 2023 as ``an Attribute in the Ethus holding two possible values instead of just one''\\
Invalid & Jan.\ 2023 on & A state that contains or triggers a contradiction, and so cannot be realized: an invalid Aethus ``will forever stay on paper, and will never exist in our physical universe'' (May 2023)\\
Child, parent & Jan.\ 2023 on & ``a child Ethus as meaning an Ethus with added information from another Ethus, then a parent Ethus would be the product of removing information from an Ethus'' (Jan.\ 2023)\\
Compression & Aug.\ 2022 on & Update without collapse: ``I like the word `compress' better'' (Aug.\ 2022); ``A superposition may never truly collapse -- rather it can only compress'' (Jan.\ 2023)\\
Rule, law, postulate & Oct.\ 2022 on & The framework's basic statements: rules (Oct.\ 2022); rules and ``basic laws,'' with its principles called ``postulates'' in a footnote (Jan.\ 2023); laws (Feb.\ 2023); postulates (May 2023 on)\\
Blank & 2024 on & Distinguished from independence: ``not statistically independent \ldots{} but it's still blank'' (draft video, June 2024)\\
Agreeing, disagreeing superposition & 2024 on & Names for the two superpositions of Section~\ref{sec:t3}, with a definition for any attribute (draft video, July 2024)\\
\bottomrule
\end{tabularx}
\end{table}

\subsection{The rules redesigned: a conditional clarification}\label{sec:t1}

The January 2023 manuscript was the first manuscript to assemble the framework's rules as three numbered rules, the first two restating rules of a five-rule list in the October 2022 math journal (``Something can only be real relative to an Ethus''; ``Any Attribute of an Anti-Ethus will be in superposition relative to its Ethus''):
\begin{quote}\small
1. Anything may only be real relative to an Ethus. 2. Any attribute of an Anti-Ethus will be in superposition relative to its Ethus. 3. An Ethus is invalid if a contradiction is triggered anywhere in the superposition of its parent Ethuses.
\end{quote}
Eleven days later by the file's metadata, on February 8, a journal entry boxed these three sentences as the ``Old Version'' and replaced them:
\begin{quote}\small
1. Anything may only be real relative to an Ethus. 2. If an Ethus is valid, then any given property will be in superposition relative to it if and only if it is statistically independent to it. 3. Any child of an invalid Ethus will also be invalid.
\end{quote}
The entry, titled as a rephrasing, says it was prompted by ``a little bit of a blend between the old second law and third law's descriptive information,'' and gives two reasons. The old third law reaches ``invalid parents based on the given Ethus working `backwards' in the tree,'' which ``makes for good practical application, but isn't very law-like''; and the old formulation ``doesn't work to hint at the clear correlation between the contradiction and the clash of Anti-Ethus superposition,'' which the new version ``fixes quite nicely.'' In the new version each concept is ``mutually exclusive,'' in the entry's phrase, from the last, and the entry singles out the new second law's contrapositive, which yields a test for invalid states, as ``a nice touch.''

Whether the redesign added or dropped content turns on entailment, and Table~\ref{tab:t1} sets out the comparison. It lists each component of the February laws with the January passage that supplies it, then each January component with its place in the February laws, and enters the interpretive assumption, if any, on which each code depends. Each direction of the second law is entered together with its contrapositive, which is one of the two cases of invalidity the entry derives from the law. Rows that carry January content through a term, such as ``superposition,'' ``child,'' and ``statistically independent,'' need no assumption beyond rule~(i): the February entry gives these terms no new sense, so it uses them in January's. Four assumptions do the work.
\begin{description}[leftmargin=1.6em, labelindent=0pt, itemsep=2pt, topsep=4pt]
\item[I1.] A contradiction triggered in a parent's superposition makes the parent invalid. January states as much (``if an Ethus with removed information from another causes a contradiction, then both Ethuses are invalid'') and asserts the resulting equivalence itself, calling the child-directed form ``another way of putting it.''
\item[I2.] To be in superposition is to hold more than one value. January says of an Anti-Ethus variable that ``it exists in all possible states at once,'' and calls an attribute of the Ethus ``holding two possible values instead of just one'' a contradiction, ``which unlike the Anti-Ethus, is strictly forbidden.''
\item[I3.] January's rules state what holds of an Ethus that can be realized, so that an Ethus that breaks one is invalid rather than a counterexample to it. January treats a breach of rule~3 this way (``rule number three would be broken if such an Ethus was permitted, and so, blatantly speaking, it is not''), and a breach of the Anti-Ethus's defining independence too: ``Even though an Anti-Ethus, in its definition, is assumed to be statistically independent of the Ethus, it still holds that any Anti-Ethus which clashes somewhere with its Ethus, hypothetical or practical, would be collectively invalid in the real sense'' (footnote~2). No passage says it of rule~2, so I3 extends the pattern.
\item[I4.] Footnote~3's check states a necessary condition for validity, not a sufficient one. Its ``Therefore'' follows from the claim that ``an invalid Ethus, in being functionally contradictory, may only beget invalid children Ethuses,'' which yields only that an Ethus with a contradicting ancestor is invalid.
\end{description}

\begin{xltabular}{\textwidth}{@{}>{\raggedright\arraybackslash}p{0.27\textwidth}>{\raggedright\arraybackslash}X>{\raggedright\arraybackslash}p{0.07\textwidth}>{\raggedright\arraybackslash}p{0.15\textwidth}@{}}
\caption{Audit of the first transition: the components of each formulation, sought in the other.}\label{tab:t1}\\
\toprule
Component & Where the other formulation supplies it & Needs & Code\\
\midrule
\endfirsthead
\toprule
Component & Where the other formulation supplies it & Needs & Code\\
\midrule
\endhead
\bottomrule
\endfoot
\multicolumn{4}{@{}p{\textwidth}@{}}{\small\emph{Construct}: the framework's conditions on states: which are real, which properties are in superposition, which states are valid. \emph{Earlier}: January rules 1--3, read with the passages of \S\S3--4 and footnotes 2--3 that state or gloss their content (``January 2023''; PDF metadata 28 Jan., self). \emph{Later}: February laws 1--3 (8 Feb.\ 2023, 4:42 PM, self; the entry quotes the January rules verbatim).}\\
\midrule
\multicolumn{4}{@{}l}{\emph{February components, sought in the January text}}\\*
\midrule
Law 1 & Rule 1, verbatim & --- & Preserved\\
Law 2, forward, with its contrapositive: an independent property of a valid Ethus is in superposition; an Ethus with an independent property out of superposition is invalid & Rule 2, with \S3: the Anti-Ethus ``contains all the properties in reality which would have zero net effect on the Ethus if they were altered. In other words, the Anti-Ethus is statistically independent to the Ethus.'' Rule 2 states this of every Ethus, which entails it of valid ones & --- & Preserved\\
Law 2, converse, with its contrapositive: a property in superposition relative to a valid Ethus is independent of it; an Ethus with a property in superposition but not independent is invalid & \S4: an attribute of the Ethus holding more than one value is ``strictly forbidden,'' and ``Either way, the Ethus in question is invalid''; \S3: a partly correlated property is ``subdivided into an Ethus component and an Anti-Ethus component.'' For a clash with the Ethus, footnotes 2--3 supply the case without I2 & I2 & Preserved on I2\\
Law 3 & \S4: ``Another way of putting it is that every child of an invalid Ethus is also invalid'' & --- & Preserved\\
\midrule
\multicolumn{4}{@{}l}{\emph{January components, sought in the February laws}}\\*
\midrule
Rule 1 & Law 1, verbatim & --- & Preserved\\
Rule 2, stated without condition & Law 2, forward, which states it of valid Ethuses; the Anti-Ethus's content carried by ``statistically independent'' & I3 & Preserved on I3; dropped if I3 is denied\\
Rule 3, parent-directed & Law 3 & I1 & Preserved on I1\\
\S4: ``no contradictions are permitted,'' a contradiction being an attribute of the Ethus holding two values & Law 2, converse, in its contrapositive, which makes such an Ethus invalid & I2 & Preserved on I2\\
\S3: the superposition ``will never arbitrarily skew one way or another without the direct order of the Ethus to do so'' & Carried by ``superposition,'' which the laws use in its January sense (rule~i; Table~\ref{tab:glossary}) & --- & Preserved\\
\S3: a partly correlated property is ``subdivided into an Ethus component and an Anti-Ethus component'' & Carried by ``statistically independent'' in its January sense (rule~i); law~2 then places each component & --- & Preserved\\
Footnote 2: an Anti-Ethus that clashes with its Ethus is invalid & Law 2, converse, in its contrapositive: footnote~3 reads the clash as a failure of independence, ``defeating its definition of independence'' & --- & Preserved\\
Footnote 3: to check an Ethus for validity, ``one must prove that every possible parent/ancestral Ethus to it will never trigger a contradiction'' & Law 3, applied along the ancestry, which gives the check's necessity & I1, I4 & Preserved on I1 and I4; its sufficiency dropped if I4 is denied\\
Child and parent, as defined in \S4 & ``Child'' used in law 3, ``parent'' as its converse; neither redefined & --- & Preserved\\
\S4: reality avoids a contradiction ``adding to the Ethus accordingly'' & Not stated or entailed in the laws & --- & Outside the construct: the update rule (see text); dropped if counted\\
\midrule
\multicolumn{4}{@{}p{\textwidth}@{}}{\small\emph{Competing reading}: denies I3. January's rule 2 holds of every Ethus, so an independent property is never out of superposition; February states the law of valid Ethuses only and makes its breach a test of invalidity, which drops the unconditional rule. \emph{Code}: clarification on I1--I4. Denying I1, I3, or I4, or counting the update rule, drops a component (contraction); denying I2 also adds the converse (revision). The Anti-Ethus coded abandonment, its work absorbed into law~2. \emph{Survival}: law~3 reappears in May 2023 as a derived consequence (Section~\ref{sec:t2}).}\\
\end{xltabular}

On I1--I4, every component of each formulation is stated or entailed in the other, and the transition is a clarification: what changed was which sentence held which content. The child-directed form moved from a gloss into a law, the Anti-Ethus left the laws with its content carried by the biconditional, and the validity condition moved into the second law. The entry's second reason describes a change of this kind, in what the laws ``hint at.'' One January statement falls outside the comparison. The manuscript's section on contradiction says that if a contradiction occurs anywhere in the superposition, ``reality will take the simplest possible path to avoid that, adding to the Ethus accordingly.'' That is a statement of how a state is updated, not of which states are real, superposed, or valid; it continues the October 2022 rule~\#5, is absent from the February laws, and returns in May (Section~\ref{sec:t2}). Counted as a component of the rules, as the competing reading of Section~\ref{sec:t2} counts it, it would be dropped, and no reading would yield a clarification.

Rule~(iii) asks for the code each denial yields. Denying I1 leaves rule~3 and the ancestral check of footnote~3 without a counterpart in the laws, a contraction. Denying I3 makes January's rule~2 a law of every Ethus, stronger than the February law, which states it of valid Ethuses only, so the unconditional rule is dropped and the transition is again a contraction. Denying I4 makes footnote~3's check sufficient as well as necessary, as the January derivation of the detector-off result uses it: ``we may remove attributes at will from our Ethus as the observer, and never once will a contradiction occur. Therefore, calling the electrons in a superposition long after the experiment is over is totally valid.'' The February laws state no sufficient condition for validity, so the sufficiency is dropped, a contraction once more. Denying I2 adds the converse direction of law~2 and leaves January's prohibition of contradiction without a counterpart in the laws, a revision. I3 and I4 are the most contestable: the first extends to rule~2 a pattern January shows elsewhere, and the second reads footnote~3 more narrowly than January uses it. Later evidence does not settle I2: the May manuscript introduces a ``converse'' of the rule as a further postulate but states the same direction both times (``any concept in superposition relative to an Aethus \ldots{} is statistically independent to it''; ``if an item is in superposition with respect to an Aethus, then it must be statistically independent to it''). My 2026 account, which calls February 8 ``the first axiomatic designation'' of the framework, is tested in Section~\ref{sec:reconstruction}.

\subsection{The third postulate: an extension}\label{sec:t2}

The February law said that invalidity passes from a state to its children. It gave no condition under which a state is protected from a contradiction lying in its future, and beyond the cases the second law's contrapositive catches, it said nothing about which states are invalid in the first place. The May 2023 manuscript builds such a condition in four steps, all on the page.

It first proposes a one-generation candidate and rejects it: ``We may think to perhaps define this as follows: `If an Aethus is valid, then no child Aethus of it may contradict any of its properties.' This is close to the final realization, but still not quite accurate.'' The same one-generation condition, in substance, stands as a ``Postulate'' in a journal entry of May 7, 2023, 10:54 PM: ``no child Aethus of mine contradicts that item in any way.'' The candidate is too strong, the manuscript explains, because a valid state can have contradicting children, since ``simply because there are other options of child Aethae to avoid that signifies that all hope was not yet lost.'' It then states the condition for a state that is, in its phrase, ``necessarily doomed'': ``For some central Aethus, if for every one of its child Aethae there exists a further child Aethus that contradicts any of its properties, then said central Aethus is invalid.'' It takes the contrapositive as the postulate:
\begin{quote}\small
Postulate 3 (The Postulate of Aethic Certainty) If some Aethus is valid, then there exists a child Aethus to it for which every one of its own child Aethae strictly do not contradict any of the properties of said Aethus.
\end{quote}
Finally it derives the earlier forms as consequences: ``every invalid Aethus must have all invalid child Aethae,'' and ``your Aethus is invalid if any one of your parent Aethae is invalid.''

Coded by components against the February laws, the transition adds two components and drops none (Table~\ref{tab:t2}). The existential--universal condition over two generations, with the criterion of invalidity that is its contrapositive, is neither stated nor entailed in those laws. Nor is the rule that a state no test shows invalid is valid (``If it also fails to prove your Aethus is invalid, then it must be valid''). On I4 the January derivation of the detector-off result relied on that rule unstated; denying I4, the rule restores the sufficiency of footnote~3's check, which the February laws dropped (Section~\ref{sec:t1}). The February third law is preserved as a derived consequence, and so, on I1, is the January rule; the second law reappears as the second postulate; ``Ethus'' becomes ``Aethus'' and ``laws'' become ``postulates,'' renamings by the sense-by-use rule (Table~\ref{tab:glossary}). The code is extension, and the candidate is coded abandonment, its rejection stated in the text that proposes it. The two-generation condition has consequences the earlier forms lack: the resolution of the January paradox (Section~\ref{sec:t3}), and a derivation that the records of different observers cannot conflict, which the January manuscript does not contain. The detector-on and detector-off predictions are not among them, since the January manuscript already derives both, the first from its third rule and the second from its second and third with the uncertainty principle and the check of footnote~3.

The strongest competing reading denies the construct criterion of Section~\ref{sec:protocol} and counts January's update rule as part of the condition. In January, reality responds to a contradiction anywhere in the superposition by ``adding to the Ethus accordingly,'' and the May text motivates the postulate by the same mechanism (``the universe diverts you away from such an invalid child Aethus by adding a property to your Aethus''). Since the addition removes the contradiction from the whole superposition, January's rule arguably has the existential--universal structure in dynamic form, which May restored. Relative to the February laws the condition is still added, so this transition's code stands, but the first transition becomes a contraction (Section~\ref{sec:t1}). On either reading the criterion of a doomed state, for which no addition can avoid a contradiction, is new in May. A weaker reading, January's check over ``every possible parent/ancestral Ethus'' turned toward children, fails on the quantifiers: a universal check over ancestors is not an existential--universal condition over descendants.

The dating is weaker than the content. The manuscript dates the two-generation form to May 2023, its deposited file produced on May 23 by its metadata. May 8 is an inference from the journal: the entry of May 7 still states the candidate and argues the detector case with it, and the entry of May 9 places ``officially cracking the measurement problem'' on the previous day, the day of the account event the arXiv message dates (Section~\ref{sec:record}). Neither entry records the form reached, so the cracking of May 8 may have completed the May 7 argument, with the two-generation form arriving by May 23. The form is independently attested, in the May wording, in the talk posted in October 2023 (43:26), and by my reading in version~1 of the preprint; the independently supported claim is existence by October 2023.

One process fact here bears on more than process. In the October 2023 talk the speaker says he ``figured I'd start with some trial and error for the third postulate, see how they perform, and then update it to make it better'' (40:17), with the double-slit cases as the test. The postulate was fitted to those cases: it accommodates the detector results rather than predicting them, and by the channel of Section~\ref{sec:evidence} those results provide no independent test of it. Accommodation is not thereby without evidential value; Hitchcock and Sober locate its characteristic risk in overfitting \cite{hitchcock2004prediction}, and the place to assess that risk here is the consequence the postulate adds beyond the fitted cases, the agreement of observers' records (Section~\ref{sec:installed}).

\begin{table}[htbp]
\centering\small
\caption{Audit of the second transition: the third postulate.}\label{tab:t2}
\begin{tabularx}{\textwidth}{@{}>{\raggedright\arraybackslash}p{0.2\textwidth}>{\raggedright\arraybackslash}X@{}}
\toprule
Construct & The framework's condition on which states can be realized\\
Earlier formulation & February law 3: ``Any child of an invalid Ethus will also be invalid,'' with law 2 (8 Feb.\ 2023, self)\\
Rejected candidate & ``If an Aethus is valid, then no child Aethus of it may contradict any of its properties,'' rejected as too strong on the page that proposes it; stated, in substance, as a postulate in the journal, 7 May 2023, 10:54 PM (self)\\
Later formulation & Postulate 3, in the May wording (quoted above), with the rule of validity by default from the May text's detector-off analysis\\
Date and basis & ``May 2023''; produced 23 May by PDF metadata; 8 May by inference from the journal entries of 7 and 9 May (self). Independently: the talk of October 2023 (43:26); preprint version 1, 14 Nov.\ 2024, by my reading\\
Added & The $\exists\forall$ condition over two generations, and the criterion of invalidity that is its contrapositive; validity by default, relied on unstated in January on I4 and otherwise the sufficiency of footnote~3's check restored\\
Preserved & Law 3 and the ancestral direction (the January rule on I1), as derived consequences; law 2, as the second postulate\\
Renamed & ``Ethus'' as ``Aethus''; ``laws'' as ``postulates''\\
Dropped & None\\
Competing reading & Denies the construct criterion of Section~\ref{sec:protocol}: January's update rule arguably has the $\exists\forall$ structure in dynamic form, restored in May; relative to February the condition is still added, and the first transition becomes a contraction. The criterion of a doomed state is new on either reading. The ancestral check turned toward children fails on the quantifiers\\
Code & Extension; the candidate coded abandonment, its rejection stated in the same text\\
Survival & The May wording recurs in the October 2023 talk and, by my reading, version 1; the 2026 statement, over ``proper'' children, keeps the quantifier structure (not coded here)\\
\bottomrule
\end{tabularx}
\end{table}

\subsection{Two superpositions: a naming and a generalization}\label{sec:t3}

The January manuscript runs the double-slit experiment with people walking through two doors on a soccer field, and ends its account with a paradox it cannot resolve. The observer at the back wall knows that each participant went through one door, but not which; ``supposing that the observer is unsure of the door the participant took seems to imply that such information would be in superposition by rule number two,'' and superposition, as the manuscript then understood it, meant going through both doors, which the third rule forbids. The footnote tries two repairs, judges that ``[b]oth of these fixes feel somewhat unsatisfactory however for their lack of simplistic clarity,'' and ends: ``Let the readers be the judge of that.''

The May manuscript draws the distinction that removes the paradox, without naming it. With the detector running, the third postulate eliminates the state in which the particle went through both slits, and what stays in superposition is that ``the particle went through one slit but not the other in both ways.'' The text anticipates the objection: ``It may seem counterintuitive how this is different from the general going through both slits at once, but to Aethic logic it is certainly distinct -- for instance, it does imply that one slit is not being passed through no matter what.'' The October 2023 talk restates the distinction, still unnamed (``A and not B plus B and not A does not equal A and B, and no nots,'' 46:42), and a summary of the same week records the want of a name: ``A clear distinction is now needed between different types of superpositions.'' The names arrive in 2024. In a draft video listed on my site under July 13, 2024, \emph{Solution to the Measurement Problem}, Part~1, the two cases are introduced as ``what I call a disagreeing superposition'' and ``what I call an agreeing superposition'' (8:21; since these are coined terms, this quotation needs checking against the audio), and, by my reading, version~1 of the preprint, dated by the archive four months later, carries both names. I date the naming itself to the turn of 2024, from recollection.

Is the 2024 name a renaming of the 2023 distinction? Component by component, partly (Table~\ref{tab:t3}). On the double-slit case the names attach to the cases the May manuscript distinguished: the disagreeing superposition of the two slits is its ``one slit but not the other in both ways,'' and the agreeing one is its ``both slits at once.'' Those components are renamed, and the insistence that the two differ is preserved. The general definition is added. The 2024 drafts define a disagreeing superposition for any attribute, one alternative at a time switched on and every other switched off, the resulting cases superposed, and an agreeing superposition as the attribute's alternatives superposed together; the 2023 texts draw the distinction only for the two slits, though the May manuscript's list of four cases for the detector-off run, ``both slits at once, the first slit but not the second, the second but not the first, or neither,'' points toward the general construction.

One component is indeterminate: the case in which the particle passes through neither slit. The May manuscript lists it among the detector-off possibilities, and its detector-on analysis reads both ways. Its statement that all that is left is for ``the particle to either go through one slit and not the other in either way'' excludes it; ``just knowing that it did not go through (at least) one of them'' and ``given that at least one slit was not traversed'' leave room for it. The October 2023 talk counts it among the detector-on alternatives (``and I guess neither,'' 46:42), the summary of the same week keeps it in both runs, and so does Part~2 of the July 13, 2024 draft (``no slits at all, the first but not the second, or the second but not the first,'' 9:05, detector on; ``particle could just miss both slits,'' 16:40, detector off). Yet both October 2023 sources also restate the detector-on result in May's two-case wording (``one slit but not the other in both ways,'' 48:26; ``to one but not the other in both ways,'' summary). Whether the May verdicts excluded the case or left it standing, the texts do not settle. On the first reading the later verdicts add a component; on the second they state one that May left implicit, and the October restatements, which keep the case while using May's wording, fit the second. On the second reading the transition is mixed, a renaming on the shared case and an extension in the general definition; on the first it is an extension throughout. Which it is turns on the indeterminate component, a finding about the evidence rather than about the transition.

The competing reading is that the May distinction was posited to absorb the January paradox, with the general typology read back into it later. The May text answers part of that reading, since there the distinction falls out of the second and third postulates applied to the detector case; Section~\ref{sec:installed} states the criterion on which that answer turns. Two later descriptions of this transition are among the claims tested in Section~\ref{sec:reconstruction}.

\begin{table}[htbp]
\centering\small
\caption{Audit of the third transition: two superpositions.}\label{tab:t3}
\begin{tabularx}{\textwidth}{@{}>{\raggedright\arraybackslash}p{0.2\textwidth}>{\raggedright\arraybackslash}X@{}}
\toprule
Construct & The distinction between two kinds of superposition\\
Problem stated & January footnote 6 (quoted above); ``January 2023'' (self)\\
Earlier formulation & ``one slit but not the other in both ways'' against ``both slits at once''; ``to Aethic logic it is certainly distinct'' (May 2023, self). Restated, unnamed, in the October 2023 talk (46:42; platform) and summary (self)\\
Later formulation & ``Disagreeing'' and ``agreeing'' superposition, defined for any attribute\\
Date and basis & Draft video \emph{Solution to the Measurement Problem}, Part 1, listed under 13 July 2024 (self; the coined terms to be checked against the audio); preprint version 1, 14 Nov.\ 2024 (archive); the naming at the turn of 2024 by recollection\\
Renamed & The two double-slit superpositions, under the new names\\
Preserved & Their difference, and its criterion: ``one slit is not being passed through no matter what''\\
Added & A definition for attributes with any number of alternatives\\
Indeterminate & The ``neither'' case: May's detector-on analysis reads both ways; kept in the October 2023 and July 2024 analyses, which also restate May's two-case result. Rival codes: added or preserved\\
Competing reading & A case-specific device installed against the January paradox and generalized in retrospect. Partly answered by the May derivation from the second and third postulates (Section~\ref{sec:installed})\\
Code & Mixed (renaming on the shared case, extension in the general definition) if May left ``neither'' implicit; extension if the later texts add it\\
\bottomrule
\end{tabularx}
\end{table}

\paragraph{The second column.} Beyond the three transitions, the record keeps what did not survive, and the protocol codes it the same way. In August 2022 the earliest design's ``unknown set'' was dropped as ``kind of an overly vague concept'' in favor of the Anti-Ethus, which left the laws in February 2023 with its work absorbed. The October 2022 rule that a contradicting superposition's ``defining Attribute \ldots{} joins its respective Ethus'' left the numbered rules in January, where invalid states took its place, but not the record: January's prose keeps its mechanism, the February laws omit it, and May uses it again (Section~\ref{sec:t2}). A ``Severed-Ethus,'' introduced on January 10, 2023 as ``the set of knowns which are permanently in the Anti-Ethus,'' appears in neither manuscript. The January footnote's two repairs were rejected on the page. ``Anti-Time,'' introduced in a journal entry of February 19, 2023 and named in the May manuscript, appears nowhere after, and the May manuscript's presentation of the first postulate as a generalization of QBism was dropped by later versions. Later still, in February 2026, the reading of quantum uncertainty as a force that conceals which-way records, shared by the 2022 and 2023 texts, was replaced by a reading on which nothing conceals anything \cite{Aethus}. My own later account calls that a change in what is fed into the postulates rather than in the postulates; the protocol would code it as a revision of the framework's physical inputs, and the postulates' own 2026 restatements separately.

% ============================================================================
\section{Testing the later accounts}\label{sec:reconstruction}

The record's later accounts are of three kinds: the narrative and historian's notes of the long manuscript \cite{Aethus}, the 2026 annotation layer of the restoration \cite{MarkedPaper}, and an earlier draft of this paper \cite{PhilSciDraft}. Each is cited in a fixed version, which the caption of Table~\ref{tab:recon} names; the dossier gives each passage's location in that version and the hash of its source file.

I selected the claims after coding the transitions. Eligible statements are those in the three sources that date the audited constructs, or describe at a date the artifacts in which they appear; where one sentence makes several such statements, each is entered as a separate claim. Among eligible statements the choice was mine and followed no rule; the tested claims are chiefly firsts, dates, and readings of early entries, one occurrence of each. Eligible statements remain untested, among them statements that place the distinction of Section~\ref{sec:t3}, or steps toward it, in 2022 and January 2023, one that finds the third postulate's mechanism in the August 2022 entry, and one that dates the May manuscript to May 8, where claim~15 tests the long manuscript's other date for it, May 20; the dossier lists those I found. The distribution of verdicts therefore describes this set, not the reliability of my retrospective accounts in general. Part of my later account was drafted within the AI-assisted collaboration described in Section~\ref{sec:methods} and adopted by me. Such passages are my account, since I adopted them, but an error in one has a different source from an error of memory, so Table~\ref{tab:recon} marks whose wording each claim is in.

{\small
\begin{xltabular}{\textwidth}{@{}>{\raggedright\arraybackslash}p{0.03\textwidth}>{\raggedright\arraybackslash}X>{\raggedright\arraybackslash}X>{\raggedright\arraybackslash}p{0.2\textwidth}@{}}
\caption{Later accounts tested against earlier artifacts, claim by claim. LM: the long manuscript, cited from the source of version 519 (PhilArchive, 19 Oct.\ 2026), a correspondence I state and the archive's version trail does not show. AL: the annotation layer, repository commit 8e54a6e (20 Oct.\ 2026). ED: an earlier draft of this paper (Sept.\ 2026). ``Mine'': in my own wording; ``drafted'': drafted within the AI-assisted collaboration and adopted by me. Grades: independent (a third party dates the evidence) or self-dated.}\label{tab:recon}\\
\toprule
& Later claim (source; wording) & Earlier evidence & Verdict (grade)\\
\midrule
\endfirsthead
\toprule
& Later claim (source; wording) & Earlier evidence & Verdict (grade)\\
\midrule
\endhead
\bottomrule
\endfoot
\multicolumn{4}{@{}l}{\emph{February 8, 2023}}\\*
1 & It was ``the first axiomatic designation'' of the framework (LM; mine) & The January manuscript states ``three fundamental rules of reality'' and calls them ``basic laws,'' having called its principles ``postulates'' in a footnote; the February entry quotes them as its ``Old Version'' & Contradicted (self-dated)\\
2 & It was ``the fundamental moment in which I first used the parent and child Aethic terms'' (LM; mine) & The January manuscript defines a child and a parent Ethus; its rule~3, and the entry's own ``Old Version,'' read ``parent Ethuses'' & Contradicted (self-dated)\\
3 & It saw ``the devising of invalid Aethae for the first time'' (LM; mine) & January rule~3 and its gloss already concern invalid states & Contradicted (self-dated)\\
\midrule
\multicolumn{4}{@{}l}{\emph{The January manuscript}}\\*
4 & It has ``the three rules assembled and numbered for the first time in manuscript'' (AL; drafted) & The October 2022 rules are numbered in the journal; no earlier manuscript states numbered rules & Supported (self-dated)\\
5 & It precedes the February entry by ``ten days'' (AL, ED; drafted) & 28 Jan.\ by PDF metadata; 8 Feb.\ by the journal & Contradicted: eleven (self-dated)\\
6 & The February entry ``boxes these exact sentences'' and ``redesigns them'' (AL; drafted) & The entry quotes the three rules verbatim as its ``Old Version'' and replaces them & Supported (self-dated)\\
7 & It boxes them ``as `the old laws' '' (AL, LM; drafted) & The box is headed ``Old Version''; the entry speaks of ``the old second law and third law'' & Contradicted as a quotation (self-dated)\\
\midrule
\multicolumn{4}{@{}l}{\emph{The January paradox}}\\*
8 & Footnote 6's paradox ``is the outstanding question the May form closes'' (LM, AL; drafted) & In May the unknown slit stays in superposition as ``one slit but not the other in both ways,'' after the third postulate eliminates ``both slits at once'' & Supported (self-dated), with a note: the distinction of Section~\ref{sec:t3} does the closing, and the third postulate makes it available\\
\midrule
\multicolumn{4}{@{}l}{\emph{The May form and its date}}\\*
9 & The first postulate has not changed since May 2023, from ``the framework in the axiomatic sense has not changed since May 8, 2023'' (LM, AL; drafted) & Stated in its May wording in the October 2023 talk (24:50); restated in new wording in the long manuscript by 2026 & Open: unchanged to Oct.\ 2023 (independent); the 2026 restatement is not coded here\\
10 & The second postulate has not changed (the same sentence; elsewhere in LM, ``content fixed, dress since evolved''; drafted) & Stated with ``statistically independent'' through July 2024 and in terms of blanks by April 2025, after the 2024 drafts distinguish the terms & Open (independent): turns on the decomposition reading (see text)\\
11 & The third postulate has not changed (the same sentence; drafted) & Stated in its May wording in the October 2023 talk (43:26) and, by my reading, version 1; restated by 2026 over ``proper'' children, in terms of validity & Open: unchanged to Oct.\ 2023 (independent), and to Nov.\ 2024 by my reading; the 2026 restatement is not coded here\\
12 & The final form dates from May 8 (the same sentence; AL: ``the two-generation form arrives May 8''; drafted) & The May 7 entry states the candidate and argues the detector case with it; the May 9 entry places ``officially cracking the measurement problem'' on the previous day; neither records the form reached & Compatible, not discriminating (self-dated): the rival is that May 8 completed the May 7 argument, and the two-generation form arrived by May 23\\
13 & The arXiv confirmation of that day dates ``my contemporaneous judgment that the work was ready'' (LM; drafted) & The message dates an account event; the May 9 entry ties the event to that judgment & Supported (self-dated link to an independently dated event)\\
14 & The May manuscript's date, May 20, has ``the same grading'' as the January manuscript's, whose day and second are given ``by the file's metadata'' (ED; drafted) & The deposited PDF's metadata gives 23 May 2023, 17:10:04 UTC & Contradicted (self-dated), reading ``the same grading'' as the same basis\\
15 & The May manuscript is ``of May 20, 2023'' (LM; drafted) & No artifact is dated May 20; the May 9 entry reports the paper already set in LaTeX; the deposited file was produced on May 23 & Open (self-dated)\\
\midrule
\multicolumn{4}{@{}l}{\emph{The candidates for the third postulate}}\\*
16 & ``the May page rejecting and refining'' a candidate as too strong (ED; drafted) & The May manuscript rejects the one-generation candidate as ``not quite accurate''; the May 7 entry had stated it, in substance, as a postulate & Supported (self-dated)\\
17 & A ``too-strong/too-weak elimination'' is ``visible across dated prints,'' the January form the too-weak candidate (ED; LM similar; drafted) & No text rejects an earlier form as too weak; the January and February forms survive as consequences; my 2026 account calls the final postulate ``slightly stronger'' & Compatible, not discriminating (self-dated): the rival is that the January form was superseded and kept as a consequence, and only the too-strong candidate was eliminated\\
\midrule
\multicolumn{4}{@{}l}{\emph{The two superpositions}}\\*
18 & Both were ``present extensionally, insisted to be distinct, months before either had a name'' (ED; LM similar; drafted) & On the double-slit case: May 2023 (self-dated) and the October 2023 talk (independent); the names first attested in July 2024 & Supported (independent, from Oct.\ 2023), with a note: ``extensionally'' holds for the double-slit cases; the general definition is the extension of Section~\ref{sec:t3}\\
19 & ``content before name'' is ``the record's recurring signature'' (ED; drafted) & One naming is coded here & Open (independent)\\
\midrule
\multicolumn{4}{@{}l}{\emph{The 2022 entries}}\\*
20 & ``statistical independence'' in the August 2022 entry was ``simply a semantic convention'' in place of blank (LM; mine) & August 2022: the unknown set dropped for independence; June 2024: ``not statistically independent \ldots{} but it's still blank'' & Open (self-dated): turns on the decomposition reading (see text)\\
21 & The push for independence moved blankness away from ``a particular pin being known or not in isolation'' (LM; mine) & The entry drops a design of ``a set of knowns and then a set of unknowns'' for independence from the Ethus, ``over the past few days'' & Supported (self-dated)\\
22 & It moved blankness to ``indistinguishability across an entire attribute of pins'' (LM; mine) & The entry speaks of ``things'' and of a ``concept''; the attribute as unit appears in the October 2022 rules: ``An Attribute may have an Anti-Ethus component and an Ethus component'' & Compatible, not discriminating (self-dated): the rival is that the entry's unit is the thing or concept, and the attribute arrives in October 2022\\
23 & The August 2022 entry ``already carries the ordering'' of probability first, the wavefunction its summary (LM; drafted) & ``the actual probability density function is very much in my Ethus, and this itself is the observed quantum wave function'' & Compatible, not discriminating (self-dated): the sentence also fits a possible conflation of the wavefunction with its associated density $|\psi|^2$\\
\end{xltabular}
}

Claims 10 and 20 need the sense-by-use rule, and they show its limits. Two readings of the 2022 and 2023 term ``statistically independent'' fit the texts. On one, the term was looser than its textbook sense: the January manuscript calls the Anti-Ethus ``statistically independent to the Ethus'' and, in the next paragraph, lets the Ethus hold ``that P(y $<$ 0) $>$ 70\%'' for an Anti-Ethus variable $y$. On the other, the texts decompose such a variable: the January manuscript divides any property ``partially correlated to the Ethus'' into ``an Ethus component and an Anti-Ethus component,'' and the October 2023 talk puts the skew in the state and leaves outside it ``a further subcomponent of the property describing the random quantile on that distribution'' (23:15), so that the term keeps its textbook sense for what remains. The 2024 drafts then separate the later pair of terms: an attribute about which a little is known is, in the speaker's words, ``not statistically independent \ldots{} but it's still blank'' (June 12, 2024, 19:38). The postulate is still stated with ``statistically independent'' in July 2024 and is described in terms of blanks by April 2025 (``with the blanks go to everything,'' 45:22 of the talk of April 27, 2025). Whether the later postulate restates the earlier one turns on the decomposition reading, which the record poses sharply without settling, so both claims are open; claim 20's ``simply a semantic convention'' also passes over the 2024 distinction. Claims 21 and 22 concern an earlier step, the replacement in August 2022 of an ``unknown set'' by a relation of independence, which the entry records within days of the change.

The first three claims are small errors, but they have the shape the protocol exists to catch. In my own history, a designation and two firsts moved to the entry that rephrased the rules, eleven days after a manuscript that already contains them; both dates are self-reported, but the February entry quotes the January rules, so the order is fixed within the record. The entry itself suggests how the move happened. Its remark that the new second law's contrapositive ``is also precisely how we initiate the concept of invalid Ethuses, which the third law after that only runs with'' concerns the order of the new laws, and read later it can pass for a chronology.

Of the twenty-three claims, the record contradicts six, supports seven, fits four without discriminating, and leaves six open. Three of the six contradictions are misdated firsts in my own wording; the other three, an arithmetic slip, a misquotation, and a misread file date, are in drafted text that I adopted. Two supported claims needed a note rather than a correction: claim~8 is right that the May form closes the January paradox, and its note says how, and claim~18's ``extensionally'' holds for the double-slit cases and not for the general definition. Claims 9 and 11 are open for a reason of their own: the first and third postulates keep their May wording in the October 2023 talk, and by 2026 both have been restated, a transition this paper does not code. None of these verdicts touches the framework's content. They concern what its author later said about when and how that content arose, which is the question the record can answer.

% ============================================================================
\section{The two-stage hypothesis, made codable}\label{sec:stages}

\begin{criterion}{Hypothesis: two stages}
Inference-driven development tends to move from formulations that rely on agent-centred analogy and familiar imagery and leave their constraints implicit, to formulations that state their constraints explicitly and remain stable across later versions; and the move is concentrated at a few transitions rather than spread evenly across the record.
\end{criterion}

I first stated this as a ``charming'' stage, in which a principle is intuited from a subtly anthropocentric vantage, giving way to a ``structural'' one. ``Charming'' cannot be coded without hindsight, so the hypothesis is restated with four markers, scored on the statement of a formulation rather than on the prose around it: (M1) reliance on agent-centred analogy, present or absent; (M2) dependence on familiar imagery, present or absent; (M3) explicitness of the constraint, where 0 means it appears only in a worked case, 1 a general statement whose key terms are fixed only by examples, and 2 a general statement whose key terms are defined in the text or in an earlier text it cites; and (M4) stability, whether the formulation recurs as a stated principle, in its wording or quantifier structure, in later versions, recurrence as a derived consequence not counting. M3 scores definability, not logical complexity. M4 is stability, not invariance under redescription, which would need equivalent descriptions of one situation, such as the January manuscript's doors and slits, compared for the verdicts a formulation gives them; that test is not run here. On a track chosen to end at a surviving form, M4 is fixed for the last formulation by the choice of track and carries no evidence there. Two historical cases suggested the hypothesis to me, Feynman's lecture ``There's Plenty of Room at the Bottom'' \cite{feynman1960room} and Einstein's route from the equivalence principle to general covariance. They are its source, not its evidence.

Three alternatives predict much the same surface. On ordinary learning, constraints become explicit gradually and formulations stabilize when they stop failing. On formalization, explicitness and stability rise while content stays fixed. On retrospective organization, later accounts impose the two-stage shape on records that do not show it. On a coded record the alternatives lead one to expect different patterns: ordinary learning, markers moving at different times; formalization, M3 rising across transitions coded clarification or renaming; retrospective organization, the shape in later accounts but not in the contemporaneous artifacts; the hypothesis, the markers moving together at a transition coded as extension or revision.

\begin{table}[htbp]
\centering\small
\caption{The markers scored on the third postulate's track.}\label{tab:markers}
\begin{tabularx}{\textwidth}{@{}>{\raggedright\arraybackslash}X>{\raggedright\arraybackslash}p{0.12\textwidth}ccc>{\raggedright\arraybackslash}p{0.22\textwidth}@{}}
\toprule
Formulation of the condition & Date & M1 & M2 & M3 & M4\\
\midrule
A future iteration of the Ethus that ``can potentially contain a current element of my Anti-Ethus'' excludes superposition & 22 Aug.\ 2022 & present & absent & 2 & No\\
Rule \#1, ``No contradictions allowed!'', applied to a superposition: ``that rule-breaking-portion of that alternative Ethus must not exist'' & 23 Oct.\ 2022 & absent & absent & 1 & Yes: restated in Jan.\ 2023 as a ``rule of reality''\\
Rule 3: invalid if a contradiction is triggered in the superposition of a parent & Jan.\ 2023 & absent & absent & 2 & No: a consequence from May, on I1\\
Law 3: the children of an invalid state are invalid & 8 Feb.\ 2023 & absent & absent & 2 & No: a consequence from May\\
Postulate 3: the $\exists\forall$ condition over two generations & May 2023 & absent & absent & 1 & Yes, by the choice of track\\
\bottomrule
\end{tabularx}
\end{table}

Scored on the third postulate's track (Table~\ref{tab:markers}), the record does not show two stages. It shows five formulations in nine months, and the markers do not move together. The agent leaves the statements of the condition by October 2022; imagery is absent from all of them, though the prose around them, May's included, still runs on agents and scenes (``your own Aethus at this moment''; ``all hope was not yet lost''). Explicitness does not rise steadily. The August 2022 condition is stated in terms the same paragraph defines (``The Anti-Ethus is simple: it is the set of all possible things which are Rebetary''), the October rule in terms fixed by examples, the January form in defined terms and the February form in terms the January manuscript defines, which the entry quotes, and the May postulate in a term, ``contradict,'' that the May text fixes by example; January defines it, but May does not cite January. M4 is Yes for the May form, as the choice of track guarantees, and for the October rule, whose prohibition January restates as ``one final important rule of reality''; the formulations between recur only as consequences.

The motive for the step to a time-independent condition is recorded twice. The October 2023 talk gives it in retrospect: the speaker ``wasn't expecting'' the framework's state ``to talk to itself from the future because that felt kind of arbitrary,'' and wanted ``a way to crack this code that was completely independent from the flow of time'' (39:30). A year earlier, the October 2022 entry had asked ``So\ldots{} are we communicating with the future???'' and answered that if so, ``it would not feel very fundamental / logically symmetric to me, so I want to find a simpler alternative to this question.'' The contemporaneous entry supports the talk's account of the motive, and shows the motive already at work a step earlier, in the October 2022 move to a separate Ethus, which became January's parent Ethuses.

Qualitatively, the track looks more like ordinary learning, with a persistent motive and markers moving at different times, than like a concentrated two-stage shift: there were more than two meaningful stages, one transition is a clarification on stated assumptions, and several constructions were abandoned on the way. That comparison is a judgment, not a measurement; a systematic comparison would need each alternative's expectations stated as scores in advance, which this paper does not attempt. The hypothesis's clause on stability cannot be tested on this track at all, since a track that ends at a surviving form is stable at its end by construction; a test needs tracks chosen without regard to survival, including tracks that end in abandonment.

The record contains genuine surprise --- the 1:42~PM continuation of the August entry opens ``Aw crap. I just thought a little bit more about this thought experiment, and it actually gets weirder still'' --- but surprise does not protect a two-stage reading from retrospective organization, since the episodes in which surprise is found are selected and classified in retrospect, by me. A genuine test needs a prediction registered before the transition it concerns, coded by rules fixed in advance, on records whose authors do not control the outcome: records other than mine.

% ============================================================================
\section{Two hazards for the coder}\label{sec:hazards}

\subsection{Installed and independently constrained structure}\label{sec:installed}

\begin{criterion}{Coding distinction: installed and independently constrained structure}
A posit is \emph{installed} when its only evidential role is introduced after an anomaly, to absorb that anomaly's residual, and it adds no independently exposed consequence. A posit is \emph{independently constrained} when the constraints from which it follows were fixed before, or independently of, the anomaly it addresses, or when it adds consequences exposed to test beyond that anomaly.
\end{criterion}

I first stated this distinction as a slogan: ``Failed theories tend to add new substance to the universe. Successful theories tend to add new reasoning to the universe.'' The slogan invites retrospective sorting, on which every successful posit becomes reasoning by being called entailed and every failed one substance by being called installed. The temporal form above can be applied to dated records before outcomes are known, and it is offered here as a candidate coding distinction requiring validation, not as a historical law. Its nearest ancestors are Einstein's distinction between principle theories and constructive theories \cite{einstein1919times}, a taxonomy of forms that makes no prediction about failure, and Lakatos's distinction between progressive and degenerating programmes \cite{lakatos1978methodology}, which reads a programme's trajectory where this distinction reads a single posit at its introduction.

Four complications keep the coding honest. Independent motivation is not deductive entailment: a posit can be motivated by constraints fixed in advance without following from them. One proposal can introduce entities, reorganize reasoning, and rescue an existing theory at once. A failed ontology can generate constraints that outlive it. And successful development often keeps some components of a proposal while rejecting others. Maxwell's electrodynamics shows all four. In 1865 he presented a theory of the field against action at a distance, citing ``[t]he mechanical difficulties \ldots{} which are involved in the assumption of particles acting at a distance,'' and assuming ``an {\ae}thereal medium filling space and permeating bodies'' \cite{maxwell1865dynamical}. The displacement current that completed his equations had entered in 1861--62 through a mechanical model of that medium, which also yielded the speed of light from electrical measurements \cite{siegel1991innovation}; the medium was later dropped and the equations kept. By the slogan, the ether is installed substance on the failed side. By the temporal form, the medium of 1861--62 was a posit with consequences exposed far beyond the anomalies it addressed, and the coding has to record both the posit's failure and the survival of what it generated. Mixed cases of this kind refine the distinction; they are not to be resolved in its favor.

Applied to the record's constructs from their introducing texts, the distinction is harder to run than its statement suggests. The Anti-Ethus was introduced in August 2022, ``over the past few days,'' to replace an ``unknown set'' its author found ``kind of an overly vague concept,'' before the human-subject puzzle the entry goes on to pose. It absorbed no anomaly, so it is not installed; but nothing in the record shows it following from constraints fixed in advance, or adding a consequence exposed beyond the cases it was built to describe. Its code is left open, and the case shows that the distinction needs a rule for posits introduced to repair a concept rather than to absorb an anomaly. The January manuscript's ``compression postulate,'' that ``[a] superposition may never truly collapse -- rather it can only compress,'' serves in footnote~6 as a rejected repair, which suggests installed; but the word arrives in August 2022 as the journal's preferred name for update, months before the paradox, so footnote~6 put an existing posit to a new use. The May distinction between the two superpositions follows from postulates fixed in the same text, which suggests independently constrained; but those postulates were fitted to the double-slit cases (Section~\ref{sec:t2}), and the one consequence they add beyond them, the agreement of observers' records, reproduces a known fact. One code is open and two are contested on the record itself, which is what a candidate distinction should show before validation; an independent coder's disagreements are the data that would validate or retire it.

\subsection{Masking}\label{sec:masking}

The sense-by-use rule answers a hazard I have called \emph{linguistic masking}. I introduced the term for one case, the description of one's own cognition in a vocabulary built around other minds \cite{MarkedPaper}. The hazard for a coder is its general form: a standing vocabulary is used for content it was not built to carry, and the words' public sense and their sense in use come apart without writer or reader noticing. It cuts in both directions in an archive. Early texts can use standard words in nonstandard senses --- whether the January manuscript's ``statistically independent'' is a case is the open question of Section~\ref{sec:reconstruction} --- and later names can import senses the earlier texts lacked, as the 2024 names would if they were read back into the 2023 distinction (Section~\ref{sec:t3}). A coder who reads early texts by textbook senses, or by the later names, will code reconstruction as continuity and continuity as reconstruction.

Masking also suggests a hypothesis about reception, stated here and not tested: readers given a novel thesis whose surface language resembles an established position, with the difference stated in the text, will tend to assimilate it to the nearest established category. A test must separate misunderstanding from disagreement, counting only paraphrases that an independent judge finds inaccurate, whatever the judge thinks of the thesis. One instance occurred while this paper was drafted, when an AI reviewer read the conditional ``if I happen to be right about my stuff'' as governing only a downstream clause, and took an evaluative claim for an evidential one; preserved in the session record, it illustrates the distinction without testing the hypothesis.

A further mechanism, by which a fierce two-sided debate entrenches what both sides share, belongs to a separate comparative study, which would need cases selected by a rule fixed in advance, cases in which the mechanism should be absent, and a specified competing explanation; a position can stay unoccupied because it is camouflaged, but also because it is unproductive, unavailable on the evidence of its day, or not distinct from positions already held \cite{stanford2006exceeding}.

% ============================================================================
\section{Judgment at the points of revision}\label{sec:art}

\begin{criterion}{The art analogy, bounded}
The record shows representations revised by judgments of form and by checks against cases, criteria that were recorded but not codified. Calling such revision an art, guided by judgment and by tips that admit exceptions, describes this recorded practice. It is not a claim that the practice is better described so than procedurally, nor that no procedure could have produced the revisions.
\end{criterion}

Four targets are easily confused when inductive practice is called an art rather than an algorithm: a universal procedure guaranteed to produce discoveries; a computational model of some discovery processes; normative guidance for assessing inductive moves; and a historical explanation of a particular discovery. The analogy denies only that the first is on offer, and claims nothing against computational models: Simon argued that scientific discovery has a logic \cite{simon1973logic}, programs built in that spirit rediscovered Kepler's third law from data \cite{langley1987scientific}, and Schmidhuber's G\"odel machines rewrite any part of their own code once they have proved the rewrite useful \cite{schmidhuber2003godel}. An argument from regress --- axiomatize inductive practice, and you write the era's assumptions into the axioms --- shows only that a method whose fixed axioms include substantive assumptions about the world cannot revise them; a fixed procedure need not contain fixed beliefs about the world. The analogy's prescriptions are tips of the kind Vonnegut gave for writing fiction, which good writers are free to break \cite{vonnegut1999bagombo}.

The analogy agrees with Norton's material theory on the negative: there are no universal inductive schemas, and ``each domain has an inductive logic native to it'' \cite{norton2021material}. It concerns a stage the material theory is not built to describe, the choice of which representation to revise when the facts that would warrant inferences within it are themselves under revision, and on that stage the record offers process evidence rather than a verdict. At three points of revision it records the criteria its author used: law-likeness, the mutual exclusivity of the laws' contents, and what the laws make visible on February 8, 2023; the rejection of two repairs ``for their lack of simplistic clarity'' in January; and in May the test of a candidate against the case it had to handle, ``close to the final realization, but still not quite accurate.'' The first two are judgments of form; the third is a check against a case. Whether such criteria are themselves governed by material facts of the domain is the open question between the two views. The protocol of Section~\ref{sec:protocol} is the complement the analysis supports: no rules of discovery in the guaranteed sense, and exacting rules of checking. In the image I have used before, the outlaw and the sheriff work different shifts on the same job.

% ============================================================================
\section{Limits, methods, and what follows}\label{sec:limits}

\subsection{Limits}\label{sec:limitsonly}

The record has one subject, who selected it, coded it, and wrote this paper. Its private leg is dated by its author. Its transitions were chosen to exercise the codes, with their outcomes known, and its later accounts were selected by the same person (Section~\ref{sec:reconstruction}). It can show how one framework was built and test what its builder later said about that; it cannot show whether the framework is true, and by itself it cannot support any generalization about discovery.

The comparison class has a selection problem of its own, and it bears on this paper's method: whose records survive to be audited is a fact about access. The notebook studies this paper sets itself beside are studies of celebrated individuals, and any appeal to the historical record of ``great individuals'' also carries an obligation stated plainly: that record is severely compromised by the systematic exclusion of women, people of color, and other marginalized groups from education and scientific institutions --- an exclusion that speaks only to the cruelty of the structures enforcing it and says nothing whatsoever about the capacities of those it targeted, the conflation of denied access with diminished ability being among the most damaging errors in the history of ideas. The irony is particularly bitter\footnote{This disconnect is one of those rare things that actually brings me physical rage: the juxtaposition of denying individuals gratitude because of race or gender with then targeting those same people for reeling under that discrimination --- for struggling to reproduce the output stats of the white men those same institutions babied.} when you also consider how the malice behind historical systemic discrimination muzzles the very class of ontology-capable individuals (people able to re-identify what a theory's basic quantities are about) we should be prizing --- while the very system doing the muzzling insists that individuals are special. A protocol like this one can be applied only to records that were kept, preserved, and made public, and it inherits the selection that decided which records exist.

\subsection{Provenance and AI assistance}\label{sec:methods}

The 2026 versions of this paper were drafted within an AI-assisted collaboration. Claude models (Anthropic) drafted connective prose and the restructuring of the analysis, and checked citations, quotations, and computations, working from my manuscripts, journal entries, and instructions; the first person, the recollections, and the decisions about what to claim are mine. The same collaboration drafted parts of the later accounts tested in Section~\ref{sec:reconstruction}, which Table~\ref{tab:recon} marks. The paper's provenance records have four statuses.
\begin{description}[leftmargin=0pt, itemsep=2pt, topsep=4pt]
\item[Completed.] The deposits listed in Table~\ref{tab:register}; the public version trail of the companion manuscript; a public ledger of the drafting sessions and a provenance record of the collaboration, in the \texttt{provenance} directory of my repository at commit 8e54a6e \cite{AethicCode}; and the supplementary dossier, which freezes every passage this paper quotes or tests, with its location and the SHA-256 hash of its source file.
\item[Planned.] Deposit of the full session record, the conversation exports with their SHA-256 hashes, which will establish the integrity of those bytes from the date of deposit, not their earlier existence; and deposit of the earlier draft of this paper, identified now by its hash in the dossier.
\item[AI-assisted checking.] Pre-submission reviews by language models, of the drafting model's family and of another family.
\item[Independent human review.] None yet; it precedes submission.
\end{description}
The logs record who contributed what to the text. They do not by themselves settle questions of intellectual authorship, which turn on what the records show about where ideas originated, the question Sections~\ref{sec:transitions} and~\ref{sec:reconstruction} take up for the framework's postulates.

\subsection{What follows}\label{sec:follows}

The analysis supports five practices, stated as practices rather than rules. Keep dated records while the work is under way, and anchor their hashes in an independent timestamp when they are written; this record's private leg was not anchored, and no later deposit can repair that. Separate what each artifact can establish before using it as evidence (Table~\ref{tab:evidence}). Code transitions component by component, from the artifacts, before consulting what survived, and state the interpretive assumptions each code needs. Test later accounts against earlier artifacts, one's own accounts first, claim by claim, and mark whose wording each is in; the errors found in Section~\ref{sec:reconstruction} would otherwise have stood in my own history. And keep abandoned constructions on the page, since a record pruned to its successes hides exactly the revisions that distinguish development from reconstruction. The single-subject limitation will dissolve only as other builders keep, and open, records of their own.

\paragraph{A note on sources.} The journal entries quoted here are from the public export of the journal \cite{AethicCode}; most are also reproduced in the long manuscript \cite{Aethus}. The 2023 manuscripts are quoted from the deposited PDFs \cite{Benander2023jan, Benander2023may}, and the summary of October 2023 from its hosted export \cite{Benander2023rules}; the talks are quoted from the transcripts described in Section~\ref{sec:record}, with timestamps for checking against the audio \cite{Benander2023talk}. The supplementary dossier gives, for every quoted or tested passage, the artifact, version, and location, and for every source file its SHA-256 hash. Recollections are mine, at most with drafted wording, and will be checkable against the session record once it is deposited.

% ============================================================================
\appendix
\section{Coding sheet}\label{app:coding}

For each transition, a coder records:
\begin{enumerate}[label=\arabic*., leftmargin=*, itemsep=1pt, topsep=4pt]
\item the construct, identified by its function in the framework;
\item the earlier formulation: passage, artifact, version, date, and dating basis (Table~\ref{tab:evidence});
\item the later formulation, with the same fields;
\item the components of each formulation, each a claim located in its text, with senses fixed by use (Table~\ref{tab:glossary});
\item for each component, a code from Table~\ref{tab:codes}, the passage in the other formulation that supplies it, and any interpretive assumption the code needs;
\item the transition's code, derived from item 5 alone, and, where an assumption is needed, the code the transition receives if the assumption is denied;
\item the strongest competing reading, stated as the denial of an assumption in item 5 or as a rival code for a named component, and the evidence that would decide between it and the code;
\item the construct's survival, recorded after item 6.
\end{enumerate}
For each later account, the coder splits the account into atomic claims, each asserting one thing, and records for each its source, version, location, and whose wording it is; the earlier artifacts consulted; and a verdict from Table~\ref{tab:codes}, with the grade of its best evidence. Where the coder's entries differ from those of Tables~\ref{tab:t1}--\ref{tab:recon}, the difference and its ground are the data the recoding exists to produce.

The primary texts are in my repository at commit 8e54a6e \cite{AethicCode}: the January and May 2023 manuscripts (\path{history/Jan2023Aethic.pdf}, \path{history/Original_Aethic.pdf}); the journal export (\path{development/journals/}, file \texttt{Foundational Journal Entries.md}); the summary of October 2023 (\path{history/MajorRulesoftheAethus.html}); and the verbatim restoration, with transcripts of the talks (\path{history/marked_original.html}). The preprint's version trail is at \url{https://philarchive.org/versions/BENARA-5}. The supplementary dossier freezes every passage used here, with its location and the SHA-256 hash of its source file.

\printbibliography

\end{document}
